%% 第31回編集委員会（物理第5巻）：組 原W2
%% p5-03.tex  原-3　電子は波でもある：ド・ブロイ波と不確定性（★3〜4・3段にしない）
%% 中身は計画書 tools/phys5-plan.md の B「原-3」。公式の持ち主は A（A23・A24・A25・A45・A46）。
%%   ほかの章は計画書 B の節の番号で指した（原-2 ③⑥⑧、原-4 ⑤、原-5 ①④、原-6、原-8 ②、原-11 ⑧、原-13・原-15）。
%% 数値はすべて Python で計算した：checks/物理第5巻/原-3/calc_p5_03.py（出力 calc_p5_03_out.txt）、
%%   code_p5_03.py（⑤の囲みのプログラムと、図の点の座標。出力 code_p5_03_out.txt・dots_p5_03.txt）。
%% 数学との重なり（N-10・W-6・Y-9）：9-4⑨（ガウスの形のはば σ と 1/σ）、4-9（うなり）、9-3（いくつもの波の和）、3-17（確率・正規分布）、
%%   8-13（波数）は指して使い、この章は「その式がどの実験と仮定から出るか」「実物の値で何が起きるか」を主にした。
%% 前の巻・数学の本との約束（計画書 A-6・C-6-2）
%%   ・Q2（力-2:909「位置と速度が決まった1つの点という考え方が使えない」）→ ①・⑦（ph:原-3:sec1・sec7）
%%   ・Q29（波-5:843「電子のような小さな粒も波としてふるまう」）→ ①②
%%   ・Q34（波-13:904「光の顕微鏡の限界 約 0.2 μm、電子顕微鏡（原-3）」）→ はてな？・⑥・例題1・12
%%   ・Q36（波-13:1635「光子を1つずつ数える干渉じま。波か粒かではなく両方」）→ ⑤
%%   ・M11（9-13:399「時間の幅と周波数の幅 → 不確定性（原-3）」）→ ⑦⑧
%%   ・M21（9-4 のガウスのはば σ と 1/σ）→ ⑦・例題10
%%   ・M19（8-13 の波数 k）→ ②⑦
%% 記号（この章）：ド・ブロイ波長 λ、運動量 p、波数 k（波-5 と同じ。カードは波-5）、ディラック定数 ħ（カード）、
%%   不確かさ Δx・Δp（カード）、角振動数 ω、群速度 v_g（波-5 は c_g。この章は粒の速さ v とくらべるので v_g。カード）、
%%   電子の波の高さ ψ（カード。原-5 で波動関数として式を立てる）、加速電圧 V、電子の質量 m（= m_e）。
%%   虚数単位はこの章では使わない（原-5 ④で i）。
%% 出典（11月11日に作業環境から確かめた）
%%   ・CODATA 2022（h・ħ・e・c・m_e・m_p・m_n）：計画書 D4 https://physics.nist.gov/cuu/Constants/Table/allascii.txt
%%   ・ノーベル物理学賞（nobelprize.org の summary）：1929 ド・ブロイ "for his discovery of the wave nature of electrons"、
%%       1932 ハイゼンベルク、1937 デビソンと G. P. トムソン "for their experimental discovery of the diffraction of electrons by crystals"、
%%       1986 ルスカ "for his fundamental work in electron optics, and for the design of the first electron microscope"、2022 量子もつれ（計画書 D37）。
%%   ・デビソンのノーベル講演（1937、https://www.nobelprize.org/uploads/2018/06/davisson-lecture.pdf）：ニッケルの結晶に電子を当てた、
%%       54 V の電子の回折、λ=h/p を確かめた、λ=12.25/√V Å の直線、1927年。
%%   ・ニッケルの原子の列の間かく 0.215 nm・50°・54 V・0.165 nm、ド・ブロイの仮説 1924年、トノムラらの実験は 50 kV（5万 V）で
%%       電子を1個ずつ、ジョンソン 1961年の電子の二重スリット、C60 の干渉 1999年：
%%       Juffmann ら "Quantum interference of molecules -- probing the wave nature of matter" arXiv:1211.3493（本文）。
%%   ・ジョンソン 1961年（Z. Phys.）・メルリら 1976年・トノムラら 1989年（Am. J. Phys. 57, 117）：Bach ら "Controlled double-slit electron
%%       diffraction" arXiv:1210.6243 の本文と文献表。
%%   ・ルスカのノーベル講演（https://www.nobelprize.org/uploads/2018/06/ruska-lecture.pdf）：1931年4月に最初の電子顕微鏡で像、
%%       1932年に 75 kV・開き角 2×10⁻² rad で分解能 2.2 Å と見こみ、約40年後に達成。
%%   ・ナトリウム 3p の遷移確率 A=0.616×10⁸ /s（589.0 nm・589.6 nm）：NIST Handbook of Basic Atomic Spectroscopic Data,
%%       https://physics.nist.gov/PhysRefData/Handbook/Tables/sodiumtable3.htm → 寿命 16.2 ns。
%%   ・ハイゼンベルクの論文 1927年（Z. Phys. 43, 172）：AIP の展示 https://history.aip.org/exhibits/heisenberg/uncertainty-relations.html
%%       （WebSearch で読んだ二次資料。ノーベル講演 heisenberg-lecture.pdf にも不確定性関係の説明がある）。
%%   ・CD のトラックの間かく 1.6 μm・DVD 0.74 μm（ECMA-130・ECMA-267）：ECMA のサイトは届かなかった。二次資料2つ
%%       （hypertextbook.com の CD の項、cdrinfo.com の DVD Technical Guide）で照らし合わせた（2026年11月）。
%%   ・熱中性子 0.0253 eV・2200 m/s は、20.4 ℃（293.6 K）の k_B T で決めた約束の値（計算で確かめた）。
%%   ・光の顕微鏡の限界 約 0.2 μm は 波-13 ⑤ の値。
%%   ・二重スリットのやってみようの表計算・プログラムの数は、code_p5_03.py を動かした結果。
\makeatletter
\tikzset{txd lab/.style={font=\sffamily\footnotesize, text=bkInk}, txd arr/.style={line width=0.9pt, -{Stealth[length=2.2mm]}},
  txd dim/.style={bkGray, line width=0.5pt, {Stealth[length=1.6mm]}-{Stealth[length=1.6mm]}},
  txd e/.style={line width=0.9pt, bkBlue!80!black, -{Stealth[length=2.2mm]}}}
% 原子（丸）
\newcommand\txdatom[1]{\filldraw[fill=bkGray!30, draw=bkInk!70, line width=0.5pt] #1 circle (1.6);}
% 電子（小さな丸と「−」）
\newcommand\txdel[1]{\filldraw[fill=bkBlue!25, draw=bkBlue!80!black, line width=0.5pt] #1 circle (1.4); \node[font=\sffamily\scriptsize, text=bkBlue!80!black] at #1 {$-$};}
% レンズ（電子のレンズ：コイルの断面）
\newcommand\txdlens[2]{\filldraw[fill=bkOrange!20, draw=bkOrange!80!black, line width=0.6pt] ($#1+(-#2,-1.6)$) rectangle ($#1+(-#2+5,1.6)$);
  \filldraw[fill=bkOrange!20, draw=bkOrange!80!black, line width=0.6pt] ($#1+(#2-5,-1.6)$) rectangle ($#1+(#2,1.6)$);}
\makeatother

\phchapter{3}{電子は波でもある：\\ド・ブロイ波と不確定性}{電子顕微鏡は、なぜふつうの顕微鏡より小さなものが見える？}{\bklevelrange{3}{4}}{電子は波でもある：ド・ブロイ波と不確定性}
\bkmichishirube{3-17,4-9,8-13,9-3,9-4,波-5,波-13,原-2}{原-4,原-5,原-6,原-8,原-13,原-15}

\begin{goal}
\begin{itemize}[label={\color{bkTheme}●}, leftmargin=1.6\zw, itemsep=0pt, topsep=1pt]
\item \bklevel{3}\enspace 加速した電圧から、電子の波長（ド・ブロイ波長）を求められる。
\item \bklevel{3}\enspace 電子の回折と、電子1個ずつの二重スリットのしま、電子顕微鏡が細かい物を見分けられるわけを、波の式で説明できる。
\item \bklevel{4}\enspace 波の集まり（波束）のはばから不確定性関係を導き、原子の中の電子のエネルギーの大きさを見積もれる。
\item \bklevel{4}\enspace 電子の波の分散関係から、群速度が粒の速さになることを示せる。
\end{itemize}
\end{goal}

\phkoushiki{%
\phk{$\lambda=\dfrac{h}{p}=\dfrac{h}{mv}$（ド・ブロイ波長）}{高校物理}{ph:原-3:debroglie}
\phk{電子線回折：$d\sin\theta=n\lambda$（$\lambda=h/p$ の電子の波で）}{高校物理}{ph:原-3:kaisetsu}
}

\begin{hatena}
\noindent\begin{minipage}[c]{0.77\linewidth}\raggedright
光の顕微鏡で見分けられるのは、どんなによいレンズを使っても約 \qty{0.2}{\micro m} までです。ウイルス（約 \qty{0.1}{\micro m}）の形はわかりません。ところが、電子顕微鏡では、原子が1つずつ並んでいるようすまで写せます。光のかわりに電子を使うと、なぜ、そんなに細かい物が見えるのでしょう。
\end{minipage}\hfill
\begin{minipage}[c]{0.2\linewidth}\centering
\begin{tikzpicture}[x=0.8mm, y=0.8mm]
  \filldraw[fill=bkGray!25, draw=bkInk, line width=0.6pt] (0,2) rectangle (10,24);
  \foreach \y in {8,14,20} {\txdlens{(5,\y)}{5}}
  \draw[bkBlue!80!black, line width=0.8pt] (5,23) -- (3.6,20) -- (6.2,14) -- (5,9) -- (5,3);
  \node[txd lab, anchor=north] at (5,1) {電子顕微鏡};
\end{tikzpicture}
\end{minipage}
\end{hatena}

\noindent 原-2 では、波だと思っていた光が、エネルギー $hf$ のつぶ（光子）としてもふるまうことを見ました。この章では反対に、小さなつぶだと思っていた電子が、波としてもふるまうことを見ます。電子の波の波長はとても短く、それが電子顕微鏡の秘密です。さらに、波であることから、電子の位置と運動量を同時にいくらでも正確に決めることはできない、という不思議なきまり（不確定性関係）が出てきます。このきまりは、原子の大きさや、原子の中の電子のエネルギーの大きさを決めています。

\begin{yattemiyou}[CD と DVD の虹をくらべよう]
\textbf{用意するもの}\par
いらなくなった CD と DVD（どちらも、文字の書いていない面を使う）、白い光の照明（部屋の電灯）\par\vspace{2pt}
\begin{enumerate}[label={\bkcircled[bkGreen]{\arabic*}}, leftmargin=2.4\zw, itemsep=1pt, topsep=0pt]
\item CD の面に部屋の電灯を映し、少しずつかたむけて、虹の色が見える向きをさがす。
\item 同じことを DVD でする。虹が見えるのは、電灯の映る向きから、どれくらい大きくかたむけたときか、CD とくらべる。
\item 青と赤では、どちらが電灯の映る向きに近いかを見る。
\end{enumerate}
\textbf{気をつけること}：太陽を映して見ない。レーザーポインターの光は使わない（目に入ると危ない）。割れた CD のふちで手を切らないように。
\end{yattemiyou}

\noindent やってみようの答え：CD の面には、細い溝（トラック）が約 \qty{1.6}{\micro m} の間かくで、DVD には約 \qty{0.74}{\micro m} の間かくで並んでいます。並んだ溝は回折格子（波-13 ③）と同じ働きをして、$d\sin\theta=\lambda$ の向きに、波長 $\lambda$ の光を強め合います。緑の光（\qty{550}{nm}）では、CD で約 20°、間かくの小さい DVD では約 48° と、大きくかたむいた向きに虹が見えます。青（波長が短い）は赤より、映る向きの近くに見えます。\textbf{規則正しく並んだ物の間かくが、波の波長と同じくらいのとき、回折で色が分かれる}のです。この章の③では、光のかわりに電子を、溝のかわりに結晶の原子の列を使った、同じ実験が出てきます。

\begin{omoidasou}
\bkrecall{原-2}{光子のエネルギーは $E=hf=hc/\lambda$（③）、運動量は $p=h/\lambda=E/c$（⑥）。$h=\qty{6.63e-34}{J.s}$、$hc=\qty{1240}{eV.nm}$。}
\bkrecall{電-2}{電荷 $e$ の電子が電圧 $V$ で加速されると、エネルギー $eV$ を得る。1 V で \qty{1}{eV}＝\qty{1.60e-19}{J}（⑥）。}
\bkrecall{波-13}{二重スリットのしまの間かく $\Delta x=L\lambda/d$（②）。回折格子 $d\sin\theta=m\lambda$（③）。見分けられる限界は波長の半分ほど（⑤）。}
\bkrecall{波-5}{進む波 $y=A\sin(\omega t-kx)$、波数 $k=2\pi/\lambda$（③）。群速度 $c_{\mathrm g}=\dd\omega/\dd k$（⑦）。}
\bkrecall{9-4}{ガウスの形 $e^{-t^2/(2\sigma^2)}$ を変換すると、周波数のはばは $1/\sigma$。時間のはば×周波数のはば＝1（⑨）。}
\bkrecall{3-17}{たくさん集めると、ばらつきは正規分布（つりがねの形）になる。はばは標準偏差 $\sigma$。}
\end{omoidasou}

%% ======================== 物理の道具 ========================
\bksection{doguu}{物理の道具}

\bksubsection{1}{つぶが波でもある：ド・ブロイの考え}
\label{ph:原-3:sec1}
第1巻『力学』では、ボールや車を「位置と速度が決まった1つの点」として扱いました。電子も、原-1 の陰極線の実験では、電場や磁場で曲がる小さなつぶとしてふるまいました。

ところが、原-2 で見たように、光は干渉や回折をする波なのに、光電効果ではエネルギー $hf$・運動量 $h/\lambda$ のつぶ（光子）としてふるまいます。1924年、フランスのド・ブロイは、これを逆に考えました。
\begin{phyomi}[ド・ブロイの考え]
「光が、波なのにつぶとしてもふるまうのなら、つぶだと思っていた電子も、波としてふるまうのではないか。光子の運動量と波長の関係 $p=h/\lambda$ は、電子にも成り立つのではないか」。
\end{phyomi}
この考えは、はじめは大胆な思いつきでした。けれども、3年後の1927年、アメリカのデビソンとガーマー、イギリスの G. P. トムソンが、結晶に電子を当てると、X線と同じような回折が起きることを見つけました（③④）。ド・ブロイは1929年に、デビソンとトムソンは1937年にノーベル物理学賞を受けました。

電子が波でもあるので、原子の中の電子は、「位置と速度が決まった1つの点」としては扱えません。第1巻（力-2）で予告したのは、このことです。どう扱えばよいかを、この章で見ていきます。

\bksubsection{2}{ド・ブロイ波長：$\lambda=h/p$}
\label{ph:原-3:sec2}
光子の運動量 $p$ と波長 $\lambda$ の関係（原-2 ⑥）を、そのまま電子などのつぶにも当てはめます。
\pheq{\lambda=\frac{h}{p}=\frac{h}{mv}}{m ＝ J·s ÷ (kg·m/s)}
\label{ph:原-3:debroglie}この波長を\textbf{ド・ブロイ波長}といい、つぶがもつ波を\textbf{物質波}（ド・ブロイ波）といいます。\bkidx{どぶろいはちょう}{ド・ブロイ波長}\bkidx{ぶっしつは}{物質波}運動量 $p$ は、質量 $m$ と速さ $v$ の積 $mv$ です（光の速さに近くないとき）。

\begin{phyomi}[$\lambda=\dfrac{h}{p}$]
「波長は、プランク定数を運動量でわったもの」と読みます。運動量が大きい（重い、または速い）つぶほど、波長は短くなります。プランク定数 $h=\qty{6.63e-34}{J.s}$ はとても小さいので、身の回りの物の波長は、ものすごく短くなります（例題4）。波としてのふるまいが目立つのは、電子のように、とても軽いつぶだけです。
\end{phyomi}

\noindent\textbf{電圧で加速した電子の波長}\par\nobreak
電子顕微鏡では、電子を電圧 $V$ で加速します。電子が得る運動エネルギーは $eV$ です（電-2 ⑥）。運動エネルギー $\tfrac12mv^2$ を運動量 $p=mv$ で書きかえると、
\[ \frac12mv^2=\frac{(mv)^2}{2m}=\frac{p^2}{2m} \]
です。これが $eV$ に等しいので、
\[ \frac{p^2}{2m}=eV\qquad\text{だから}\qquad p=\sqrt{2meV} \]
これを $\lambda=h/p$ に入れて、
\pheq{\lambda=\frac{h}{\sqrt{2meV}}}{m ＝ J·s ÷ $\sqrt{\text{kg·C·V}}$}
\label{ph:原-3:kasoku}です。電子の質量 $m=\qty{9.11e-31}{kg}$、$e=\qty{1.60e-19}{C}$ を入れると、
\[ \lambda=\frac{1.226}{\sqrt{V}}\ \mathrm{nm}\qquad(V\ \text{はボルトの数}) \]
となります。デビソンも、1937年の講演で、同じ形の式（0.1 nm を単位にして $12.25/\sqrt V$）と実験の値をくらべています。

\begin{bkcenter}
\begin{bktable}{colspec={Q[l,m,wd=31mm] *{6}{Q[c,m,wd=13.5mm]}}, colsep=3pt, rowsep=1pt, row{1}={bg=bkTheme!10}}
加速電圧 $V$ & 1 V & 100 V & 1 kV & 10 kV & 100 kV & 200 kV \\
$\lambda=h/\sqrt{2meV}$〔pm〕 & 1226 & 123 & 38.8 & 12.3 & 3.88 & 2.74 \\
相対論で直した波長〔pm〕 & 1226 & 123 & 38.8 & 12.2 & 3.70 & 2.51 \\
\end{bktable}
\bktabcaption{加速電圧と電子の波長（1 pm＝$10^{-12}$ m＝0.001 nm）。100 V で原子の間かく（約 0.1〜0.3 nm）と同じくらい。数十 kV をこえると、相対論の直しが要る（原-11 ⑧）。}\label{ph:原-3:tab-lambda}
\end{bkcenter}

\begin{phbaisuru}[4倍]{加速電圧 $V$}
\noindent\begin{minipage}[c]{0.47\linewidth}\raggedright
$\lambda$ は $\sqrt V$ に反比例するので、電圧を4倍にすると、波長は $1/\sqrt4=1/2$ になります。100 V で 0.123 nm の電子は、400 V で 0.061 nm です。電圧を100倍にすると、波長は $1/10$ です。両方の目もりを対数にしたグラフ（3-12）では、傾き $-\tfrac12$ の直線になります。点線は相対論で直した波長で、数十 kV から少しずつ短くなります。
\end{minipage}\hfill
\begin{minipage}[c]{0.5\linewidth}\centering
\begin{tikzpicture}
\begin{axis}[bkplot, axis lines=left, width=56mm, height=46mm, xmode=log, ymode=log, xmin=1, xmax=3e5, ymin=1e-3, ymax=2, xtick={1,100,1e4}, xticklabels={1,100,\num{10000}}, ytick={0.001,0.01,0.1,1}, yticklabels={0.001,0.01,0.1,1}, xlabel={$V$（V）}, ylabel={$\lambda$（nm）}]
  \addplot[bkcurve, domain=1:3e5, samples=40] {1.2264/sqrt(x)};
  \addplot[bkcurve2, densely dotted, domain=1e3:3e5, samples=40] {1.2264/sqrt(x*(1+x/1021998))};
\end{axis}
\end{tikzpicture}
\end{minipage}
\end{phbaisuru}

\begin{phflow}[光子の運動量から、電子の波長へ]
\phflowrow{$p=\dfrac{h}{\lambda}$（光子）}{つぶにも\\同じ関係}{$\lambda=\dfrac{h}{p}=\dfrac{h}{mv}$}
\phflowrow{$\lambda=\dfrac{h}{p}$}{電圧 $V$ で加速\\$\dfrac{p^2}{2m}=eV$}{$\lambda=\dfrac{h}{\sqrt{2meV}}$}
\end{phflow}

\begin{seikaku}
電子が光の速さに近づくと、運動量は $mv$ ではなくなります（相対論。原-11 ⑧で導く）。その結果だけを使うと、$p=\sqrt{2meV\,(1+eV/(2mc^2))}$ です。$mc^2$ は電子の質量をエネルギーにした \qty{511}{keV}（原-11）なので、10 kV では1\,\%、100 kV では約10\,\%、運動量の2乗が大きくなり、波長はその平方根の分（100 kV で約5\,\%）短くなります。電子顕微鏡の 100〜300 kV では、この直しが欠かせません。
\end{seikaku}

\begin{phmodel}[電子の平面波]
\phignore{電子どうしの力。加速する前の電子のエネルギー（0 とする）。}
\phassume{電子は決まった運動量 $p$ で進み、波長 $h/p$ の波がどこまでも続く（平面波）。光の速さよりずっと遅い。}
\phlimit{電子のいる場所を小さくしぼるとき（波束：⑦）。数十 kV より高い電圧（相対論）。}
\end{phmodel}

\Needspace{12\baselineskip}\bksubsection{3}{結晶で電子が回折する：デビソンとガーマーの実験}
\label{ph:原-3:sec3}
\noindent\begin{minipage}[c]{0.45\linewidth}\raggedright
アメリカのベル研究所のデビソンとガーマーは、ニッケルの結晶の表面に、54 V で加速した電子を、表面に垂直に当てました。はね返る電子の数を、向き（角 $\theta$）を変えてはかると、$\theta=\ang{50}$ の向きに、はね返る電子がとくに多くなりました（1927年）。

結晶の表面には、原子が間かく $a=\qty{0.215}{nm}$ で列になって並んでいます。となりどうしの原子ではね返った波の道のりの差は $a\sin\theta$ です（図\ref{ph:原-3:fig-dg}）。
\end{minipage}\hfill
\begin{minipage}[c]{0.53\linewidth}\centering
\begin{tikzpicture}[x=1mm, y=1mm]
  \foreach \x in {0,12,24,36,48} {\txdatom{(\x,0)}}
  \fill[bkGray!12] (-4,-1.6) rectangle (52,-5); \node[txd lab, anchor=north] at (46,-5.3) {結晶};
  \draw[txd e] (12,26) -- (12,3.4);
  \draw[txd e] (24,26) -- (24,3.4);
  \draw[bkRed!80!black, line width=0.9pt, -{Stealth[length=2.2mm]}] (12,2) -- ++(40:24);
  \draw[bkRed!80!black, line width=0.9pt, -{Stealth[length=2.2mm]}] (24,2) -- ++(40:24);
  \draw[bkInk, line width=0.5pt] (12,12) arc (90:40:10);
  \node[txd lab] at (15.6,14.6) {$\theta$};
  \draw[bkInk, line width=0.4pt] (24,2) -- ($(12,2)+(40:{12*cos(40)})$);
  \draw[bkOrange!85!black, line width=1.6pt] (12,2) -- ($(12,2)+(40:{12*cos(40)})$);
  \node[txd lab, text=bkOrange!85!black, anchor=east] at (10.6,6.5) {差 $a\sin\theta$};
  \draw[txd dim] (36,4.2) -- (48,4.2); \node[txd lab, anchor=south] at (42,4.6) {$a$};
  \node[txd lab, anchor=south] at (24,27) {54 V の電子};
  \node[txd lab, anchor=west] at (42,22) {はね返る電子};
\end{tikzpicture}
\bkfigcaption{ニッケルの表面の原子の列ではね返る電子。}\label{ph:原-3:fig-dg}
\end{minipage}
\par\vspace{3pt}
道のりの差が波長のちょうど整数倍なら、すべての列ではね返った波の山がそろって強め合います。CD の溝と同じ、反射型の回折格子の式（波-13 ③）です。
\pheq{a\sin\theta=n\lambda\qquad(n=1,\ 2,\ \cdots)}{m ＝ m}
\label{ph:原-3:kaisetsu}54 V の電子のド・ブロイ波長と、$n=1$ の向きは
\[ \lambda=\frac{1.226}{\sqrt{54}}\ \mathrm{nm}=\qty{0.167}{nm},\qquad \sin\theta=\frac{0.167}{0.215}=0.776,\qquad\theta\approx\ang{51} \]
となり、実験の 50° とよく合いました（例題3）。反対に、実験の 50° から波長を求めると $0.215\times\sin\ang{50}=\qty{0.165}{nm}$ で、ド・ブロイの式の値と1\,\%ほどしか違いません。デビソンたちは、電圧をいろいろ変えて、同じことを確かめました。

\begin{phmijika}[CD の虹と、同じしくみ]
やってみようの CD では、間かく \qty{1.6}{\micro m} の溝が、波長 \qty{550}{nm} の光を約 20° に強め合いました。ニッケルの表面では、間かく \qty{0.215}{nm} の原子の列が、波長 \qty{0.167}{nm} の電子の波を約 50° に強め合います。物の大きさは1万分の1になりましたが、「間かく」と「波長」の比が同じくらいなので、同じような角になるのです。
\end{phmijika}

\begin{phbaisuru}{加速電圧 $V$}
波長が $1/\sqrt2$ になり、$\sin\theta$ も $1/\sqrt2$ 倍になるので、強め合う向き $\theta$ は小さくなります（表面に垂直な向きに近づく）。54 V で約 51° だった向きは、108 V では約 33° です。電圧を上げると、$n=2$（道のりの差が2波長）の向きも現れます（$\sin\theta=2\lambda/a\leqq1$ になるのは、約 130 V より上）。
\end{phbaisuru}

\bksubsection{4}{電子線回折：うすい膜を通すとリングになる}
\label{ph:原-3:sec4}
同じ1927年、G. P. トムソンは、もっと高い電圧（数万 V）で加速した電子を、金属のうすい膜に通しました。膜の向こうに置いた写真の板には、同心円のリングがいくつも写りました（図\ref{ph:原-3:fig-ring}）。これを\textbf{電子線回折}といいます。\bkidx{でんしせんかいせつ}{電子線回折}

\noindent\begin{minipage}[c]{0.5\linewidth}\raggedright
うすい膜は、小さな結晶のつぶが、いろいろな向きを向いて集まったものです。結晶の中の原子の面（間かく $d$）で強め合う向きは、ブラッグの条件 $2d\sin\theta=n\lambda$ で決まります（結晶の面での回折は、原-8 ②で X線とともに導きます）。どの向きのつぶでも同じ角で強め合うので、強め合った電子は、まっすぐ進む向きの周りに円すいの形に広がり、板の上ではリングになります（図\ref{ph:原-3:fig-cone}）。

角が小さいとき、リングの半径は、膜から板までの距離を $L$ とすると、およそ $R\approx L\lambda/d$ です。\textbf{リングが大きいほど、波長が長いか、面の間かくがせまい}のです。
\end{minipage}\hfill
\begin{minipage}[c]{0.47\linewidth}\centering
\begin{tikzpicture}[x=1mm, y=1mm]
  \foreach \X/\k/\t in {0/1/{電圧 $V$}, 31/0.7071/{電圧 $2V$}}{
    \fill[bkInk!85] (\X-14,-14) rectangle (\X+14,14);
    \fill[white] (\X,0) circle (1.0);
    \foreach \r/\w in {4.5/1.2, 5.2/0.8, 7.4/0.9, 8.6/0.7, 9.1/0.5, 10.4/0.5, 11.4/0.4, 11.7/0.4}
      {\draw[white, line width=\w pt] (\X,0) circle ({\r*\k});}
    \node[txd lab, anchor=north] at (\X,-15) {\t};}
\end{tikzpicture}
\bkfigcaption{金属のうすい膜の電子線回折（形をかいた図。写真の板は黒い地に白いリング）。電圧を2倍にすると、どのリングも $1/\sqrt2$ 倍に小さくなる。}\label{ph:原-3:fig-ring}
\end{minipage}

\begin{phbaisuru}{加速電圧 $V$}
波長が $1/\sqrt2$ になるので、リングの半径もすべて $1/\sqrt2$ 倍に小さくなります。リングの半径の比（どのリングがどれの何倍か）は変わりません。比は結晶の原子の並び方で決まるからです。だから、リングの並びを見れば、電圧を知らなくても、どんな結晶かがわかります。
\end{phbaisuru}

\begin{bkfigure}
\begin{tikzpicture}[x=1mm, y=1mm]
  \draw[txd e] (0,0) -- (22,0);
  \node[txd lab, anchor=south] at (10,1) {電子};
  \fill[bkGray!35] (23,-14) rectangle (24.4,14);
  \node[txd lab, anchor=north] at (23.7,-15) {うすい膜};
  \draw[bkGray, line width=0.5pt, dashed] (24,0) -- (104,0);
  \draw[bkRed!80!black, line width=0.9pt, -{Stealth[length=2mm]}] (24,0) -- (104,16);
  \draw[bkRed!80!black, line width=0.9pt, -{Stealth[length=2mm]}] (24,0) -- (104,-16);
  \filldraw[fill=bkInk!12, draw=bkInk, line width=0.6pt] (104,-22) rectangle (107,22);
  \node[txd lab, anchor=south] at (105.5,23) {写真の板};
  \fill[bkRed!80!black] (105.5,16) circle (0.9) (105.5,-16) circle (0.9);
  \draw[txd dim] (110,0) -- (110,16); \node[txd lab, anchor=west] at (111,8) {$R$};
  \draw[txd dim] (24,-26) -- (104,-26); \node[txd lab, anchor=north] at (64,-27) {$L$};
  \draw[bkInk, line width=0.5pt] (44,0) arc (0:11.3:20); \node[txd lab, anchor=west] at (45,2.2) {$2\theta$};
  \node[txd lab, anchor=south west, align=left] at (30,17) {結晶の面で強め合った電子は、\\まっすぐの向きから $2\theta$ 曲がり、\\円すいの形に広がる};
\end{tikzpicture}
\bkfigcaption{電子線回折でリングができるわけ（横から見た図）。いろいろな向きの小さな結晶のつぶのうち、ブラッグの条件をみたすつぶで強め合った電子が、まっすぐの向きの周りに円すいの形に広がり、板の上で半径 $R$ のリングになる。$\theta$ が小さいとき $R\approx L\times2\theta\approx L\lambda/d$。}\label{ph:原-3:fig-cone}
\end{bkfigure}

電子線回折は、いまでは材料を調べる大切な道具です。電子顕微鏡の中で、見たい場所だけに電子を当てて回折のリングや点の並びを写すと、そこにある結晶の種類と向きがわかります（工学ではこう使う）。

\Needspace{12\baselineskip}\bksubsection{5}{電子の二重スリット：1個ずつ届いて、しまになる}
\label{ph:原-3:sec5}
波であることのいちばんの証拠は、二重スリットの干渉（波-13 ②）です。電子でも、1961年にドイツのジョンソンが、ごく細いすきまを2つ並べて、電子のしまを写しました。

では、電子を\textbf{1個ずつ}送ったらどうなるでしょう。1989年、日本の外村彰らのグループは、電子顕微鏡の中で電子を 50 kV で加速し、1個がスクリーンに届いてから次の1個が出るほど弱い電子の流れで、届いた位置を1個ずつ記録しました（すきまのかわりに、電子の流れを2つに分けて重ねる「電子線バイプリズム」という道具を使いました）。

\begin{bkfigure}
\begin{tikzpicture}[x=1mm, y=1mm]
  \fill[bkBlue!80!black] (-6.2,4.3) circle (0.4) (0.3,-1.9) circle (0.4) (16.7,10.3) circle (0.4) (-7.4,-11.0) circle (0.4) (1.7,-12.7) circle (0.4) (-10.3,-2.1) circle (0.4) (-1.6,8.5) circle (0.4) (7.0,2.1) circle (0.4) (0.9,6.6) circle (0.4) (9.0,-0.3) circle (0.4) (9.3,10.8) circle (0.4) (5.7,5.8) circle (0.4)
    (11.3,-3.3) circle (0.4) (5.9,-1.1) circle (0.4) (-11.6,-12.7) circle (0.4) (9.4,-7.3) circle (0.4) (-7.8,-4.1) circle (0.4) (-12.2,-7.7) circle (0.4) (-9.3,2.6) circle (0.4) (-8.3,-2.3) circle (0.4);
  \fill[bkBlue!80!black] (44.5,5.1) circle (0.35) (54.4,-2.9) circle (0.35) (28.7,-8.4) circle (0.35) (32.2,-12.9) circle (0.35) (26.5,-11.5) circle (0.35) (43.4,-11.4) circle (0.35) (55.5,-13.7) circle (0.35) (50.9,-3.8) circle (0.35) (45.1,-4.9) circle (0.35) (52.4,0.5) circle (0.35) (35.6,13.2) circle (0.35) (34.4,-8.5) circle (0.35)
    (52.8,-9.3) circle (0.35) (54.1,8.7) circle (0.35) (51.9,-10.8) circle (0.35) (54.7,9.4) circle (0.35) (31.1,13.6) circle (0.35) (33.9,-6.0) circle (0.35) (34.9,-6.9) circle (0.35) (41.2,11.8) circle (0.35) (44.1,0.7) circle (0.35) (32.9,8.6) circle (0.35) (51.9,-5.0) circle (0.35) (44.6,8.2) circle (0.35)
    (35.2,7.8) circle (0.35) (44.3,7.5) circle (0.35) (44.4,-1.4) circle (0.35) (45.2,12.7) circle (0.35) (55.9,-10.9) circle (0.35) (42.0,-7.5) circle (0.35) (44.5,-8.8) circle (0.35) (41.9,-4.8) circle (0.35) (36.0,-2.2) circle (0.35) (35.5,12.1) circle (0.35) (42.4,1.5) circle (0.35) (34.8,-9.7) circle (0.35)
    (53.1,0.0) circle (0.35) (32.7,-13.9) circle (0.35) (35.3,6.7) circle (0.35) (45.8,-0.7) circle (0.35) (55.2,-0.1) circle (0.35) (44.2,13.9) circle (0.35) (35.4,-12.0) circle (0.35) (35.0,-1.2) circle (0.35) (26.5,-0.7) circle (0.35) (44.1,0.1) circle (0.35) (26.7,-11.8) circle (0.35) (40.5,-13.8) circle (0.35)
    (37.1,2.5) circle (0.35) (45.0,4.5) circle (0.35) (61.0,6.5) circle (0.35) (54.9,0.7) circle (0.35) (44.2,8.8) circle (0.35) (50.8,-13.5) circle (0.35) (43.6,8.6) circle (0.35) (52.7,1.0) circle (0.35) (49.6,6.8) circle (0.35) (35.3,-6.7) circle (0.35) (52.0,6.9) circle (0.35) (35.4,-5.6) circle (0.35)
    (46.4,-12.6) circle (0.35) (35.9,5.5) circle (0.35) (44.6,-1.0) circle (0.35) (59.3,-1.2) circle (0.35) (42.2,-8.4) circle (0.35) (46.9,0.7) circle (0.35) (59.8,9.2) circle (0.35) (44.3,5.9) circle (0.35) (43.5,-14.3) circle (0.35) (43.7,-5.7) circle (0.35) (55.9,12.3) circle (0.35) (36.6,-3.1) circle (0.35)
    (36.1,-11.5) circle (0.35) (55.7,12.6) circle (0.35) (35.2,0.3) circle (0.35) (33.1,13.1) circle (0.35) (51.7,6.7) circle (0.35) (35.6,-5.7) circle (0.35) (46.0,-9.6) circle (0.35) (32.2,11.7) circle (0.35) (43.9,11.7) circle (0.35) (35.5,11.2) circle (0.35) (52.7,-2.5) circle (0.35) (44.8,-4.7) circle (0.35)
    (60.4,0.1) circle (0.35) (34.1,-7.2) circle (0.35) (57.1,-13.5) circle (0.35) (43.1,-14.4) circle (0.35) (60.5,-11.3) circle (0.35) (53.3,1.5) circle (0.35) (44.9,-3.2) circle (0.35) (34.3,-14.1) circle (0.35) (34.7,13.3) circle (0.35) (51.2,-13.8) circle (0.35) (43.9,-2.3) circle (0.35) (35.5,12.2) circle (0.35)
    (34.4,-4.7) circle (0.35) (52.4,-8.5) circle (0.35) (55.2,-8.0) circle (0.35) (53.1,13.0) circle (0.35) (43.9,-8.0) circle (0.35) (60.6,-12.9) circle (0.35) (52.6,4.7) circle (0.35) (26.6,-4.3) circle (0.35) (59.9,13.4) circle (0.35) (33.8,9.3) circle (0.35) (55.3,-13.0) circle (0.35) (43.1,12.1) circle (0.35)
    (33.3,11.4) circle (0.35) (27.9,-12.6) circle (0.35) (35.7,-5.3) circle (0.35) (36.1,7.4) circle (0.35) (59.5,-7.7) circle (0.35) (35.3,-0.2) circle (0.35) (53.9,-8.7) circle (0.35) (52.9,-12.5) circle (0.35) (35.8,-11.6) circle (0.35) (43.9,-1.5) circle (0.35) (34.7,3.5) circle (0.35) (33.5,-7.3) circle (0.35)
    (46.7,-3.0) circle (0.35) (34.6,4.4) circle (0.35) (61.2,-0.7) circle (0.35) (36.8,-9.0) circle (0.35) (59.6,2.8) circle (0.35) (33.6,2.9) circle (0.35) (26.9,5.1) circle (0.35) (33.0,-8.5) circle (0.35) (54.3,-12.6) circle (0.35) (45.8,-11.8) circle (0.35) (52.0,-14.2) circle (0.35) (55.2,11.0) circle (0.35)
    (42.6,-14.0) circle (0.35) (45.8,11.8) circle (0.35) (44.7,-11.3) circle (0.35) (43.7,13.4) circle (0.35) (33.4,12.8) circle (0.35) (34.3,10.0) circle (0.35) (55.5,-8.1) circle (0.35) (35.1,11.4) circle (0.35) (53.0,9.7) circle (0.35) (45.8,-5.6) circle (0.35) (42.1,-13.7) circle (0.35) (34.7,8.1) circle (0.35)
    (42.5,-0.8) circle (0.35) (41.6,-1.7) circle (0.35) (44.4,3.9) circle (0.35) (53.7,-12.8) circle (0.35) (44.1,13.0) circle (0.35) (55.0,13.9) circle (0.35) (43.7,12.0) circle (0.35) (35.7,-5.2) circle (0.35) (27.8,-9.8) circle (0.35) (54.0,0.9) circle (0.35) (54.4,14.1) circle (0.35) (46.7,7.6) circle (0.35)
    (42.0,7.0) circle (0.35) (35.4,13.9) circle (0.35) (37.4,-13.4) circle (0.35) (34.5,-13.8) circle (0.35) (26.6,-11.3) circle (0.35) (33.6,-0.7) circle (0.35) (35.0,-11.6) circle (0.35) (53.9,-6.8) circle (0.35) (46.2,4.2) circle (0.35) (52.2,11.6) circle (0.35) (33.5,0.2) circle (0.35) (32.6,-5.8) circle (0.35)
    (53.9,-14.2) circle (0.35) (42.8,-1.4) circle (0.35) (34.4,-7.7) circle (0.35) (52.7,9.9) circle (0.35) (51.4,1.6) circle (0.35) (44.8,4.1) circle (0.35) (60.3,10.9) circle (0.35) (27.0,-7.6) circle (0.35) (52.6,10.9) circle (0.35) (51.9,-5.5) circle (0.35) (33.9,-12.2) circle (0.35) (27.4,12.4) circle (0.35)
    (27.5,5.5) circle (0.35) (52.3,2.6) circle (0.35) (53.0,-5.2) circle (0.35) (41.3,-7.0) circle (0.35) (45.3,10.4) circle (0.35) (32.5,-8.6) circle (0.35) (43.7,-9.1) circle (0.35) (43.8,9.6) circle (0.35) (36.4,5.7) circle (0.35) (43.9,3.9) circle (0.35) (33.7,11.5) circle (0.35) (44.0,11.1) circle (0.35)
    (34.7,0.9) circle (0.35) (52.9,4.2) circle (0.35) (53.6,-10.8) circle (0.35) (34.8,-5.7) circle (0.35) (31.9,-7.3) circle (0.35) (32.1,-9.0) circle (0.35) (52.2,-8.8) circle (0.35) (33.7,-13.4) circle (0.35);
  \fill[bkBlue!80!black] (85.2,6.4) circle (0.24) (86.4,3.7) circle (0.24) (97.9,3.7) circle (0.24) (79.2,-1.3) circle (0.24) (97.7,-0.3) circle (0.24) (104.6,1.2) circle (0.24) (89.4,0.4) circle (0.24) (88.8,12.9) circle (0.24) (77.8,-3.1) circle (0.24) (97.2,13.9) circle (0.24) (80.1,-14.0) circle (0.24) (94.9,-6.8) circle (0.24)
    (78.4,12.7) circle (0.24) (88.9,8.7) circle (0.24) (86.9,-7.9) circle (0.24) (99.1,-5.9) circle (0.24) (89.7,9.6) circle (0.24) (79.9,-7.1) circle (0.24) (85.4,-14.3) circle (0.24) (95.8,-7.3) circle (0.24) (78.7,-4.4) circle (0.24) (98.2,11.3) circle (0.24) (89.1,-1.8) circle (0.24) (79.8,-10.4) circle (0.24)
    (87.8,-0.9) circle (0.24) (87.9,10.4) circle (0.24) (86.9,4.8) circle (0.24) (87.5,-13.4) circle (0.24) (88.9,-8.3) circle (0.24) (85.7,1.6) circle (0.24) (99.4,9.4) circle (0.24) (80.0,13.7) circle (0.24) (70.7,12.1) circle (0.24) (77.9,-1.2) circle (0.24) (97.2,-9.2) circle (0.24) (98.0,-7.0) circle (0.24)
    (71.9,8.9) circle (0.24) (86.1,-3.2) circle (0.24) (84.9,4.2) circle (0.24) (77.9,-13.9) circle (0.24) (75.0,6.3) circle (0.24) (79.0,-9.0) circle (0.24) (96.0,3.7) circle (0.24) (95.3,12.5) circle (0.24) (95.6,-14.0) circle (0.24) (96.0,10.4) circle (0.24) (87.5,-3.8) circle (0.24) (94.2,8.6) circle (0.24)
    (90.3,-12.6) circle (0.24) (99.5,3.0) circle (0.24) (98.8,-14.4) circle (0.24) (79.7,2.5) circle (0.24) (90.5,10.1) circle (0.24) (98.4,7.2) circle (0.24) (86.8,-7.1) circle (0.24) (96.8,-1.2) circle (0.24) (97.5,2.3) circle (0.24) (88.8,2.6) circle (0.24) (77.1,-9.2) circle (0.24) (95.0,1.8) circle (0.24)
    (75.7,14.3) circle (0.24) (71.2,14.1) circle (0.24) (90.3,8.0) circle (0.24) (97.3,13.3) circle (0.24) (79.4,-8.6) circle (0.24) (76.8,-12.9) circle (0.24) (79.5,4.1) circle (0.24) (105.2,-13.3) circle (0.24) (79.5,11.6) circle (0.24) (105.0,-10.2) circle (0.24) (79.5,-0.5) circle (0.24) (76.4,-10.3) circle (0.24)
    (94.2,6.3) circle (0.24) (77.3,12.3) circle (0.24) (86.3,9.3) circle (0.24) (87.2,-10.3) circle (0.24) (78.3,6.1) circle (0.24) (89.9,10.7) circle (0.24) (79.8,-9.9) circle (0.24) (87.7,11.6) circle (0.24) (87.2,-7.0) circle (0.24) (77.6,14.1) circle (0.24) (95.9,13.8) circle (0.24) (88.9,-1.9) circle (0.24)
    (90.5,-9.2) circle (0.24) (97.5,-8.7) circle (0.24) (80.1,3.2) circle (0.24) (90.9,-12.5) circle (0.24) (96.1,6.4) circle (0.24) (87.2,-6.7) circle (0.24) (86.1,-10.9) circle (0.24) (95.4,7.2) circle (0.24) (88.5,1.1) circle (0.24) (78.3,-11.4) circle (0.24) (71.5,5.7) circle (0.24) (77.3,2.9) circle (0.24)
    (95.1,10.6) circle (0.24) (95.6,-10.8) circle (0.24) (87.8,-6.3) circle (0.24) (100.6,2.1) circle (0.24) (96.6,9.4) circle (0.24) (88.9,12.7) circle (0.24) (97.7,-7.5) circle (0.24) (100.2,0.5) circle (0.24) (79.2,3.8) circle (0.24) (88.2,-10.4) circle (0.24) (101.5,4.1) circle (0.24) (79.8,-6.5) circle (0.24)
    (88.7,13.0) circle (0.24) (80.6,4.2) circle (0.24) (79.9,1.0) circle (0.24) (75.7,-10.6) circle (0.24) (80.8,-6.1) circle (0.24) (79.0,1.3) circle (0.24) (86.6,3.2) circle (0.24) (86.9,-7.4) circle (0.24) (78.3,-3.4) circle (0.24) (91.0,-4.2) circle (0.24) (90.0,-6.2) circle (0.24) (105.1,7.8) circle (0.24)
    (76.0,10.7) circle (0.24) (85.9,-3.2) circle (0.24) (96.4,7.5) circle (0.24) (87.1,6.0) circle (0.24) (101.0,7.7) circle (0.24) (86.5,-6.0) circle (0.24) (88.0,-9.1) circle (0.24) (81.1,2.2) circle (0.24) (90.9,12.1) circle (0.24) (77.7,11.8) circle (0.24) (70.9,1.9) circle (0.24) (87.9,7.9) circle (0.24)
    (88.6,5.3) circle (0.24) (88.9,-3.6) circle (0.24) (105.4,0.9) circle (0.24) (99.1,-5.2) circle (0.24) (88.4,11.4) circle (0.24) (85.2,-6.0) circle (0.24) (88.4,-9.0) circle (0.24) (98.1,5.5) circle (0.24) (70.6,9.9) circle (0.24) (77.4,1.5) circle (0.24) (79.8,0.9) circle (0.24) (84.9,-9.9) circle (0.24)
    (97.1,-11.5) circle (0.24) (98.7,-14.0) circle (0.24) (97.5,6.2) circle (0.24) (79.8,11.6) circle (0.24) (86.2,-12.8) circle (0.24) (78.3,10.8) circle (0.24) (87.5,-7.4) circle (0.24) (76.6,-9.0) circle (0.24) (90.2,1.0) circle (0.24) (99.4,-7.3) circle (0.24) (77.2,-7.6) circle (0.24) (95.2,-6.6) circle (0.24)
    (99.7,-1.6) circle (0.24) (95.8,13.2) circle (0.24) (88.2,-1.4) circle (0.24) (79.1,-8.3) circle (0.24) (88.5,-6.6) circle (0.24) (97.5,4.5) circle (0.24) (90.4,-3.2) circle (0.24) (90.2,0.0) circle (0.24) (80.9,-5.4) circle (0.24) (78.4,6.2) circle (0.24) (80.4,11.8) circle (0.24) (100.7,-6.4) circle (0.24)
    (95.0,10.0) circle (0.24) (76.9,11.9) circle (0.24) (96.2,-13.2) circle (0.24) (71.9,-8.7) circle (0.24) (81.1,-13.3) circle (0.24) (79.4,5.3) circle (0.24) (82.7,9.6) circle (0.24) (97.7,-13.2) circle (0.24) (96.7,5.6) circle (0.24) (99.5,-11.8) circle (0.24) (97.2,5.8) circle (0.24) (98.7,1.3) circle (0.24)
    (89.5,-12.7) circle (0.24) (77.6,-10.5) circle (0.24) (78.8,0.4) circle (0.24) (90.2,9.1) circle (0.24) (86.2,-5.1) circle (0.24) (86.9,-9.6) circle (0.24) (73.0,14.2) circle (0.24) (96.8,6.3) circle (0.24) (85.7,1.2) circle (0.24) (79.3,-10.4) circle (0.24) (99.9,-9.0) circle (0.24) (96.5,-9.1) circle (0.24)
    (99.4,-3.8) circle (0.24) (89.8,9.6) circle (0.24) (78.9,1.9) circle (0.24) (88.0,-5.8) circle (0.24) (90.8,5.4) circle (0.24) (104.5,-13.2) circle (0.24) (85.9,6.4) circle (0.24) (85.4,4.7) circle (0.24) (88.5,-4.7) circle (0.24) (76.3,-1.6) circle (0.24) (90.2,-8.8) circle (0.24) (87.3,-5.8) circle (0.24)
    (71.5,5.1) circle (0.24) (86.2,4.6) circle (0.24) (90.3,-11.1) circle (0.24) (89.7,10.4) circle (0.24) (79.4,0.9) circle (0.24) (79.3,1.6) circle (0.24) (78.4,-11.1) circle (0.24) (88.5,-12.1) circle (0.24) (84.8,-1.8) circle (0.24) (95.8,9.5) circle (0.24) (97.6,7.9) circle (0.24) (96.6,7.6) circle (0.24)
    (95.8,3.0) circle (0.24) (98.6,12.2) circle (0.24) (77.3,9.7) circle (0.24) (90.9,-13.8) circle (0.24) (77.0,-6.0) circle (0.24) (89.3,8.9) circle (0.24) (103.7,-4.5) circle (0.24) (71.7,9.2) circle (0.24) (87.2,9.9) circle (0.24) (72.0,-4.0) circle (0.24) (86.9,-3.2) circle (0.24) (82.9,2.3) circle (0.24)
    (82.2,-1.2) circle (0.24) (105.0,-10.2) circle (0.24) (80.1,-6.9) circle (0.24) (97.6,3.9) circle (0.24) (96.4,3.9) circle (0.24) (84.7,-4.7) circle (0.24) (87.4,-7.4) circle (0.24) (78.7,-10.5) circle (0.24) (86.4,2.0) circle (0.24) (81.1,-12.5) circle (0.24) (81.1,6.5) circle (0.24) (90.9,-9.3) circle (0.24)
    (79.5,-9.7) circle (0.24) (78.1,-8.6) circle (0.24) (91.5,8.9) circle (0.24) (77.9,7.2) circle (0.24) (86.7,12.2) circle (0.24) (95.8,-11.7) circle (0.24) (76.0,7.7) circle (0.24) (79.5,7.6) circle (0.24) (97.7,7.8) circle (0.24) (80.2,12.7) circle (0.24) (75.1,-0.7) circle (0.24) (78.5,-11.8) circle (0.24)
    (71.5,-12.7) circle (0.24) (88.1,-9.2) circle (0.24) (75.7,6.8) circle (0.24) (88.0,-4.5) circle (0.24) (98.5,-5.1) circle (0.24) (96.2,4.2) circle (0.24) (78.3,-6.8) circle (0.24) (85.7,-8.5) circle (0.24) (97.1,-5.8) circle (0.24) (105.3,2.0) circle (0.24) (80.4,-0.4) circle (0.24) (86.4,11.0) circle (0.24)
    (80.9,-11.2) circle (0.24) (104.6,11.9) circle (0.24) (94.4,-4.6) circle (0.24) (71.7,8.4) circle (0.24) (94.6,0.1) circle (0.24) (78.6,-11.4) circle (0.24) (88.0,-4.2) circle (0.24) (87.9,-4.3) circle (0.24) (80.1,-6.8) circle (0.24) (72.9,0.3) circle (0.24) (89.5,5.5) circle (0.24) (95.6,-5.2) circle (0.24)
    (84.8,14.4) circle (0.24) (78.9,-11.1) circle (0.24) (94.8,4.2) circle (0.24) (89.7,13.6) circle (0.24) (78.7,7.0) circle (0.24) (89.5,3.2) circle (0.24) (87.1,0.6) circle (0.24) (87.2,-10.3) circle (0.24) (88.8,9.0) circle (0.24) (78.8,9.3) circle (0.24) (86.6,9.4) circle (0.24) (79.3,-4.1) circle (0.24)
    (98.6,-1.4) circle (0.24) (86.2,8.6) circle (0.24) (97.1,5.2) circle (0.24) (78.8,-4.6) circle (0.24) (83.8,-8.6) circle (0.24) (90.5,-9.3) circle (0.24) (95.6,-5.1) circle (0.24) (80.8,4.3) circle (0.24) (98.9,-3.3) circle (0.24) (95.4,4.8) circle (0.24) (82.0,7.4) circle (0.24) (87.9,7.4) circle (0.24)
    (86.7,9.8) circle (0.24) (87.7,9.4) circle (0.24) (84.6,5.2) circle (0.24) (97.5,-8.0) circle (0.24) (96.9,-12.8) circle (0.24) (87.4,5.5) circle (0.24) (80.1,-5.4) circle (0.24) (89.9,-2.3) circle (0.24) (77.6,1.1) circle (0.24) (87.6,-3.9) circle (0.24) (98.6,12.1) circle (0.24) (90.0,-5.4) circle (0.24)
    (89.2,1.9) circle (0.24) (98.4,6.1) circle (0.24) (98.6,-1.3) circle (0.24) (101.2,-1.9) circle (0.24) (100.7,2.8) circle (0.24) (87.9,5.0) circle (0.24) (80.8,9.7) circle (0.24) (77.7,-11.7) circle (0.24) (94.5,-12.8) circle (0.24) (94.8,9.7) circle (0.24) (80.8,2.4) circle (0.24) (89.6,-1.0) circle (0.24)
    (96.0,5.2) circle (0.24) (88.9,-4.4) circle (0.24) (96.8,12.4) circle (0.24) (89.4,-10.4) circle (0.24) (100.1,12.2) circle (0.24) (86.2,7.0) circle (0.24) (80.9,1.7) circle (0.24) (98.1,-1.1) circle (0.24) (75.7,10.9) circle (0.24) (89.3,-0.9) circle (0.24) (88.2,-8.7) circle (0.24) (101.7,12.8) circle (0.24)
    (87.6,2.0) circle (0.24) (94.6,7.2) circle (0.24) (75.9,-13.5) circle (0.24) (78.7,8.2) circle (0.24) (95.4,-7.3) circle (0.24) (96.2,-10.1) circle (0.24) (77.4,-4.9) circle (0.24) (98.2,-0.1) circle (0.24) (96.8,-8.7) circle (0.24) (98.0,6.3) circle (0.24) (96.6,9.5) circle (0.24) (98.6,-1.9) circle (0.24)
    (103.6,13.5) circle (0.24) (98.1,9.7) circle (0.24) (78.6,-1.2) circle (0.24) (78.8,11.8) circle (0.24) (98.3,12.7) circle (0.24) (99.4,-11.9) circle (0.24) (71.1,9.0) circle (0.24) (86.5,-8.3) circle (0.24) (76.9,13.2) circle (0.24) (88.0,0.2) circle (0.24) (105.1,9.5) circle (0.24) (70.5,12.8) circle (0.24)
    (79.3,-11.5) circle (0.24) (88.5,12.3) circle (0.24) (88.8,11.7) circle (0.24) (87.9,13.7) circle (0.24) (96.8,-5.0) circle (0.24) (96.7,-6.3) circle (0.24) (97.2,1.3) circle (0.24) (94.7,-9.7) circle (0.24) (98.3,-10.4) circle (0.24) (88.6,-9.7) circle (0.24) (96.3,-5.4) circle (0.24) (104.2,-5.5) circle (0.24)
    (81.5,7.4) circle (0.24) (87.0,4.9) circle (0.24) (96.8,-4.0) circle (0.24) (90.3,12.7) circle (0.24) (79.8,-12.4) circle (0.24) (98.8,-5.6) circle (0.24) (97.6,12.7) circle (0.24) (96.5,-3.1) circle (0.24) (98.7,-9.0) circle (0.24) (88.7,11.6) circle (0.24) (88.1,4.8) circle (0.24) (90.6,9.9) circle (0.24)
    (96.9,11.5) circle (0.24) (99.3,-8.3) circle (0.24) (78.3,5.1) circle (0.24) (81.9,3.7) circle (0.24) (98.8,-13.5) circle (0.24) (97.7,0.1) circle (0.24) (85.3,3.7) circle (0.24) (95.9,-2.9) circle (0.24) (88.4,-7.9) circle (0.24) (86.4,-8.9) circle (0.24) (88.6,-8.7) circle (0.24) (95.1,-6.7) circle (0.24)
    (79.7,-9.7) circle (0.24) (81.2,-9.4) circle (0.24) (98.7,-11.4) circle (0.24) (97.7,14.1) circle (0.24) (87.2,-10.7) circle (0.24) (88.3,-7.2) circle (0.24) (96.1,9.0) circle (0.24) (100.0,-14.0) circle (0.24) (78.0,-3.5) circle (0.24) (98.0,-13.2) circle (0.24) (87.5,-8.5) circle (0.24) (91.2,-8.6) circle (0.24)
    (86.5,3.9) circle (0.24) (95.2,-12.5) circle (0.24) (95.9,-0.8) circle (0.24) (95.5,4.3) circle (0.24) (91.5,-7.5) circle (0.24) (105.0,-3.1) circle (0.24) (71.8,13.7) circle (0.24) (98.6,-13.0) circle (0.24) (78.2,12.4) circle (0.24) (79.7,-13.9) circle (0.24) (97.0,13.6) circle (0.24) (95.3,8.8) circle (0.24)
    (88.1,2.7) circle (0.24) (98.5,-12.8) circle (0.24) (89.6,-3.0) circle (0.24) (98.9,-10.9) circle (0.24) (85.5,13.7) circle (0.24) (89.6,-11.7) circle (0.24) (100.2,-3.8) circle (0.24) (70.7,-11.1) circle (0.24) (88.7,-2.7) circle (0.24) (95.2,-5.7) circle (0.24) (96.8,-7.2) circle (0.24) (84.9,-7.8) circle (0.24)
    (98.3,-11.4) circle (0.24) (99.4,3.6) circle (0.24) (104.3,0.9) circle (0.24) (104.4,-8.3) circle (0.24) (71.0,6.2) circle (0.24) (105.1,-1.8) circle (0.24) (96.6,-7.2) circle (0.24) (79.5,5.5) circle (0.24) (77.8,13.4) circle (0.24) (94.8,-12.6) circle (0.24) (72.8,-4.0) circle (0.24) (75.5,-0.2) circle (0.24)
    (76.7,-5.7) circle (0.24) (80.2,-8.9) circle (0.24) (75.5,13.9) circle (0.24) (90.4,-12.5) circle (0.24) (79.3,0.7) circle (0.24) (98.4,-3.4) circle (0.24) (97.3,-8.8) circle (0.24) (78.2,-3.9) circle (0.24) (100.9,4.7) circle (0.24) (99.8,-13.6) circle (0.24) (85.8,-1.2) circle (0.24) (95.4,-11.0) circle (0.24)
    (87.3,-7.8) circle (0.24) (85.9,-12.5) circle (0.24) (96.5,10.7) circle (0.24) (88.9,-9.5) circle (0.24) (96.1,-1.3) circle (0.24) (81.6,4.7) circle (0.24) (76.8,4.4) circle (0.24) (88.5,-0.9) circle (0.24) (96.4,-7.7) circle (0.24) (88.0,4.3) circle (0.24) (78.4,11.6) circle (0.24) (96.7,12.6) circle (0.24)
    (78.7,0.3) circle (0.24) (86.5,13.3) circle (0.24) (79.5,1.8) circle (0.24) (95.4,-7.5) circle (0.24) (75.7,0.6) circle (0.24) (86.3,11.2) circle (0.24) (80.6,13.4) circle (0.24) (89.0,8.8) circle (0.24) (104.2,-12.2) circle (0.24) (79.4,-3.8) circle (0.24) (88.8,-12.3) circle (0.24) (96.5,-0.8) circle (0.24)
    (78.7,13.2) circle (0.24) (77.7,13.3) circle (0.24) (96.2,4.7) circle (0.24) (97.3,-8.0) circle (0.24) (78.0,-13.4) circle (0.24) (85.2,2.4) circle (0.24) (85.8,-0.5) circle (0.24) (104.3,7.7) circle (0.24) (81.0,9.3) circle (0.24) (77.4,1.0) circle (0.24) (95.9,-12.4) circle (0.24) (70.7,5.6) circle (0.24)
    (70.6,-6.8) circle (0.24) (71.2,12.5) circle (0.24) (104.5,-4.8) circle (0.24) (88.8,-2.4) circle (0.24) (87.3,-12.8) circle (0.24) (94.9,8.4) circle (0.24) (96.0,-0.2) circle (0.24) (90.1,-10.0) circle (0.24) (95.6,8.6) circle (0.24) (98.6,2.5) circle (0.24) (78.6,4.7) circle (0.24) (90.3,-8.9) circle (0.24)
    (90.9,3.9) circle (0.24) (78.9,-2.2) circle (0.24) (89.2,-13.5) circle (0.24) (95.9,-6.0) circle (0.24) (77.7,4.7) circle (0.24) (79.6,-7.9) circle (0.24) (98.3,-5.3) circle (0.24) (95.7,3.1) circle (0.24) (88.5,12.4) circle (0.24) (77.4,0.2) circle (0.24) (88.8,6.2) circle (0.24) (89.0,-11.4) circle (0.24)
    (87.4,4.8) circle (0.24) (78.3,-0.7) circle (0.24) (86.4,-6.4) circle (0.24) (86.8,-5.8) circle (0.24) (97.5,-11.9) circle (0.24) (84.3,13.3) circle (0.24) (97.4,-11.3) circle (0.24) (89.1,-9.5) circle (0.24) (85.4,-10.8) circle (0.24) (80.6,6.2) circle (0.24) (96.2,-8.8) circle (0.24) (104.0,9.0) circle (0.24)
    (79.4,14.1) circle (0.24) (81.7,7.1) circle (0.24) (79.8,-3.6) circle (0.24) (96.9,12.1) circle (0.24) (99.0,4.5) circle (0.24) (96.6,4.4) circle (0.24) (89.7,-4.6) circle (0.24) (97.9,10.7) circle (0.24) (78.0,-7.2) circle (0.24) (78.4,-12.5) circle (0.24) (96.4,-1.1) circle (0.24) (96.2,10.8) circle (0.24)
    (90.6,2.5) circle (0.24) (87.4,11.0) circle (0.24) (85.5,0.7) circle (0.24) (81.1,6.7) circle (0.24) (77.0,0.0) circle (0.24) (98.6,3.0) circle (0.24) (79.6,-2.3) circle (0.24) (88.5,-11.7) circle (0.24) (70.7,6.2) circle (0.24) (96.7,-7.0) circle (0.24) (88.6,3.0) circle (0.24) (96.7,6.1) circle (0.24)
    (98.4,-10.3) circle (0.24) (86.6,14.3) circle (0.24) (79.1,-9.8) circle (0.24) (81.1,-8.1) circle (0.24) (76.9,-1.7) circle (0.24) (77.4,-10.2) circle (0.24) (76.4,5.8) circle (0.24) (96.0,3.7) circle (0.24) (85.6,1.6) circle (0.24) (98.2,5.1) circle (0.24) (89.7,-9.1) circle (0.24) (98.5,-5.1) circle (0.24)
    (87.3,-3.9) circle (0.24) (79.6,-3.2) circle (0.24) (89.4,-4.8) circle (0.24) (99.3,5.5) circle (0.24) (71.3,-12.7) circle (0.24) (98.7,-7.8) circle (0.24) (85.7,7.2) circle (0.24) (96.9,13.7) circle (0.24) (71.8,-10.8) circle (0.24) (76.2,3.1) circle (0.24) (85.7,-7.1) circle (0.24) (76.6,-10.0) circle (0.24)
    (94.3,-1.8) circle (0.24) (98.2,-1.4) circle (0.24) (75.5,-3.4) circle (0.24) (94.7,8.7) circle (0.24) (98.0,-14.2) circle (0.24) (98.4,4.9) circle (0.24) (79.7,-9.8) circle (0.24) (97.0,7.5) circle (0.24) (97.7,-6.4) circle (0.24) (79.5,4.3) circle (0.24) (95.8,4.4) circle (0.24) (80.2,12.1) circle (0.24)
    (87.1,7.2) circle (0.24) (89.4,-10.6) circle (0.24) (86.8,6.3) circle (0.24) (96.0,12.5) circle (0.24) (77.0,9.1) circle (0.24) (88.0,-9.7) circle (0.24) (98.0,11.9) circle (0.24) (97.5,1.0) circle (0.24) (96.6,-1.1) circle (0.24) (84.7,-10.4) circle (0.24) (87.0,-6.7) circle (0.24) (97.2,-9.7) circle (0.24)
    (89.7,-13.6) circle (0.24) (77.7,-6.2) circle (0.24) (72.8,1.5) circle (0.24) (77.7,-11.0) circle (0.24) (95.5,-13.0) circle (0.24) (99.4,-8.6) circle (0.24) (88.5,13.4) circle (0.24) (104.8,7.1) circle (0.24) (80.0,-3.9) circle (0.24) (89.8,11.5) circle (0.24) (79.4,-9.8) circle (0.24) (98.5,0.1) circle (0.24)
    (71.0,-2.2) circle (0.24) (78.8,-7.4) circle (0.24) (99.3,-11.7) circle (0.24) (87.2,-4.7) circle (0.24) (97.3,4.9) circle (0.24) (88.1,3.0) circle (0.24) (76.8,-12.0) circle (0.24) (82.1,4.3) circle (0.24) (85.3,0.4) circle (0.24) (95.2,-9.6) circle (0.24) (72.1,10.0) circle (0.24) (88.2,-11.4) circle (0.24)
    (96.5,0.3) circle (0.24) (70.6,-2.9) circle (0.24) (77.4,5.4) circle (0.24) (105.3,-6.7) circle (0.24) (85.6,-12.4) circle (0.24) (88.1,-8.8) circle (0.24) (98.8,12.5) circle (0.24) (99.0,-10.3) circle (0.24) (87.4,5.4) circle (0.24) (104.7,6.2) circle (0.24) (98.5,-5.1) circle (0.24) (95.8,5.6) circle (0.24)
    (85.8,-11.8) circle (0.24) (94.0,-13.2) circle (0.24) (105.4,11.8) circle (0.24) (90.3,0.5) circle (0.24) (87.4,-3.6) circle (0.24) (88.4,-8.0) circle (0.24) (80.2,0.1) circle (0.24) (77.0,-6.5) circle (0.24) (97.5,8.1) circle (0.24) (88.6,12.3) circle (0.24) (88.4,-6.0) circle (0.24) (99.1,6.6) circle (0.24)
    (78.0,-4.1) circle (0.24) (96.9,-8.4) circle (0.24) (78.7,4.5) circle (0.24) (94.8,-12.6) circle (0.24) (98.7,-4.7) circle (0.24) (96.3,-3.8) circle (0.24) (80.0,6.1) circle (0.24) (85.2,13.7) circle (0.24) (76.9,-6.3) circle (0.24) (79.4,-3.2) circle (0.24) (83.8,-3.7) circle (0.24) (97.1,5.2) circle (0.24)
    (81.8,-5.0) circle (0.24) (88.3,-8.3) circle (0.24) (79.0,-3.5) circle (0.24) (98.9,7.9) circle (0.24) (78.5,8.0) circle (0.24) (96.8,14.2) circle (0.24) (77.9,-11.7) circle (0.24) (87.8,10.4) circle (0.24) (80.4,13.1) circle (0.24) (98.2,9.6) circle (0.24) (71.5,13.3) circle (0.24) (87.9,-8.7) circle (0.24)
    (94.8,4.2) circle (0.24) (89.6,8.1) circle (0.24) (90.5,-1.7) circle (0.24) (96.9,-1.5) circle (0.24) (91.9,12.5) circle (0.24) (88.7,11.2) circle (0.24) (76.1,-6.4) circle (0.24) (78.5,12.2) circle (0.24) (89.8,10.4) circle (0.24) (86.7,-1.3) circle (0.24) (99.6,5.6) circle (0.24) (88.1,7.3) circle (0.24)
    (88.9,12.0) circle (0.24) (87.0,8.3) circle (0.24) (78.1,-14.4) circle (0.24) (98.4,9.5) circle (0.24) (89.8,3.6) circle (0.24) (98.0,9.1) circle (0.24) (88.3,-8.0) circle (0.24) (75.7,2.2) circle (0.24) (77.2,1.1) circle (0.24) (85.9,9.7) circle (0.24) (77.5,7.7) circle (0.24) (78.1,-5.6) circle (0.24)
    (95.9,2.1) circle (0.24) (87.3,4.1) circle (0.24) (94.7,1.5) circle (0.24) (96.2,9.8) circle (0.24) (101.0,-7.2) circle (0.24) (79.3,-1.1) circle (0.24) (95.8,-8.0) circle (0.24) (70.8,-11.3) circle (0.24) (105.4,-13.3) circle (0.24) (97.0,12.4) circle (0.24) (88.4,-1.8) circle (0.24) (87.2,-2.7) circle (0.24)
    (86.7,-7.3) circle (0.24) (87.0,-5.8) circle (0.24) (77.7,4.5) circle (0.24) (93.5,-1.7) circle (0.24) (98.2,1.3) circle (0.24) (98.8,6.6) circle (0.24) (85.0,0.6) circle (0.24) (99.6,7.3) circle (0.24) (80.3,-13.2) circle (0.24) (77.6,7.6) circle (0.24) (80.2,0.6) circle (0.24) (90.1,-13.4) circle (0.24)
    (95.4,12.2) circle (0.24) (95.6,-3.6) circle (0.24) (71.1,-8.0) circle (0.24) (77.3,-10.5) circle (0.24) (99.2,-7.2) circle (0.24) (100.7,-12.6) circle (0.24) (95.7,-0.1) circle (0.24) (77.0,-6.3) circle (0.24) (80.1,-11.9) circle (0.24) (98.5,-7.8) circle (0.24) (88.1,-3.6) circle (0.24) (90.0,10.5) circle (0.24)
    (76.6,3.3) circle (0.24) (77.3,6.1) circle (0.24) (85.4,5.6) circle (0.24) (96.4,-2.8) circle (0.24) (99.9,7.0) circle (0.24) (87.1,-3.8) circle (0.24) (76.5,2.0) circle (0.24) (77.9,-11.4) circle (0.24) (81.1,-3.0) circle (0.24) (83.8,10.4) circle (0.24) (83.6,13.2) circle (0.24) (78.5,-13.7) circle (0.24)
    (79.1,13.5) circle (0.24) (100.2,6.1) circle (0.24) (79.8,-8.9) circle (0.24) (97.8,2.9) circle (0.24) (90.3,12.3) circle (0.24) (97.7,-1.1) circle (0.24) (97.1,-11.2) circle (0.24) (79.2,11.3) circle (0.24) (87.8,12.9) circle (0.24) (98.3,-9.4) circle (0.24) (104.5,-12.2) circle (0.24) (86.3,-0.6) circle (0.24)
    (97.8,-7.6) circle (0.24) (88.2,3.1) circle (0.24) (105.4,-4.7) circle (0.24) (70.9,-12.8) circle (0.24) (90.8,-3.4) circle (0.24) (79.7,-12.5) circle (0.24) (89.5,-0.2) circle (0.24) (84.7,6.7) circle (0.24) (98.5,-11.1) circle (0.24) (90.7,-5.9) circle (0.24) (97.1,-4.8) circle (0.24) (87.2,1.0) circle (0.24)
    (71.0,10.8) circle (0.24) (98.5,-3.4) circle (0.24) (78.4,0.8) circle (0.24) (79.1,-3.3) circle (0.24) (71.1,-0.2) circle (0.24) (76.0,-5.7) circle (0.24) (88.3,9.2) circle (0.24) (99.2,12.6) circle (0.24) (98.9,-14.4) circle (0.24) (97.6,9.6) circle (0.24) (105.0,9.6) circle (0.24) (80.1,1.2) circle (0.24)
    (97.3,-11.6) circle (0.24) (86.8,-9.8) circle (0.24) (79.0,6.2) circle (0.24) (98.0,5.7) circle (0.24) (80.8,6.9) circle (0.24) (78.0,-6.3) circle (0.24) (87.3,12.1) circle (0.24) (79.5,-11.4) circle (0.24) (88.1,-9.9) circle (0.24) (87.3,9.8) circle (0.24) (88.5,-1.6) circle (0.24) (99.2,9.5) circle (0.24)
    (78.0,-0.9) circle (0.24) (84.9,8.6) circle (0.24) (96.1,13.8) circle (0.24) (78.3,-2.4) circle (0.24) (85.4,-5.3) circle (0.24) (87.8,13.8) circle (0.24) (90.2,-12.9) circle (0.24) (78.8,-1.7) circle (0.24) (89.3,-13.2) circle (0.24) (81.1,1.5) circle (0.24) (86.1,-9.0) circle (0.24) (86.4,-4.4) circle (0.24)
    (98.7,4.4) circle (0.24) (100.0,-1.6) circle (0.24) (95.7,-7.7) circle (0.24) (95.2,6.4) circle (0.24) (105.0,10.9) circle (0.24) (81.5,8.2) circle (0.24) (87.6,7.8) circle (0.24) (78.1,-14.4) circle (0.24) (79.6,7.6) circle (0.24) (96.1,-8.5) circle (0.24) (89.2,6.6) circle (0.24) (77.4,-4.5) circle (0.24)
    (96.0,-9.4) circle (0.24) (96.1,7.2) circle (0.24) (75.7,-7.2) circle (0.24) (88.5,11.1) circle (0.24) (96.8,1.8) circle (0.24) (104.0,-4.9) circle (0.24) (81.7,-2.0) circle (0.24) (99.8,-14.0) circle (0.24) (77.9,-3.8) circle (0.24) (79.2,-1.6) circle (0.24) (97.6,-4.5) circle (0.24) (90.2,-4.9) circle (0.24)
    (78.8,9.6) circle (0.24) (104.9,6.1) circle (0.24) (89.0,8.5) circle (0.24) (100.6,-8.1) circle (0.24) (97.3,12.5) circle (0.24) (77.5,12.7) circle (0.24) (97.6,2.9) circle (0.24) (95.8,5.3) circle (0.24) (87.0,-5.7) circle (0.24) (89.6,-0.3) circle (0.24) (104.0,5.9) circle (0.24) (96.0,-13.9) circle (0.24)
    (80.1,12.2) circle (0.24) (80.0,-10.5) circle (0.24) (77.8,11.4) circle (0.24) (103.6,-2.3) circle (0.24) (87.8,-10.2) circle (0.24) (88.4,-10.6) circle (0.24) (79.8,-1.2) circle (0.24) (78.0,-4.9) circle (0.24) (81.5,-2.4) circle (0.24) (87.5,5.7) circle (0.24) (77.7,2.7) circle (0.24) (97.6,-4.5) circle (0.24)
    (104.7,-8.4) circle (0.24) (96.6,8.0) circle (0.24) (88.6,12.1) circle (0.24) (97.0,-10.7) circle (0.24) (85.4,13.6) circle (0.24) (95.8,-1.5) circle (0.24) (78.9,10.8) circle (0.24) (86.8,-12.7) circle (0.24) (78.9,5.7) circle (0.24) (98.4,-3.2) circle (0.24) (78.7,10.5) circle (0.24) (100.5,2.7) circle (0.24)
    (86.6,-12.6) circle (0.24) (96.1,8.9) circle (0.24) (96.1,-7.8) circle (0.24) (99.0,6.3) circle (0.24) (71.7,1.8) circle (0.24) (76.5,3.4) circle (0.24) (79.4,-6.5) circle (0.24) (97.9,-6.9) circle (0.24) (76.3,-10.6) circle (0.24) (78.0,-3.0) circle (0.24) (97.1,-0.7) circle (0.24) (80.0,11.1) circle (0.24)
    (86.6,-8.4) circle (0.24) (76.7,5.2) circle (0.24) (89.5,6.1) circle (0.24) (85.7,-6.9) circle (0.24) (89.8,-7.2) circle (0.24) (80.7,-4.4) circle (0.24) (72.1,-7.0) circle (0.24) (89.0,2.5) circle (0.24) (85.5,13.9) circle (0.24) (96.9,-13.1) circle (0.24) (104.9,10.0) circle (0.24) (80.9,-8.7) circle (0.24)
    (80.4,-11.9) circle (0.24) (105.0,3.2) circle (0.24) (79.2,3.7) circle (0.24) (79.0,-0.7) circle (0.24) (97.9,-5.9) circle (0.24) (98.3,-9.9) circle (0.24) (84.6,-3.3) circle (0.24) (89.8,-2.8) circle (0.24) (98.5,-4.9) circle (0.24) (89.2,2.1) circle (0.24) (89.3,5.9) circle (0.24) (98.9,-5.2) circle (0.24)
    (88.5,11.0) circle (0.24) (97.1,-4.0) circle (0.24) (98.5,-12.0) circle (0.24) (89.0,1.1) circle (0.24) (103.1,-14.1) circle (0.24) (78.8,-10.0) circle (0.24) (89.1,-9.9) circle (0.24) (96.7,-0.6) circle (0.24) (97.9,-5.6) circle (0.24) (96.1,8.6) circle (0.24) (95.5,10.1) circle (0.24) (90.5,-6.9) circle (0.24)
    (86.2,11.8) circle (0.24) (86.1,3.0) circle (0.24) (99.4,-10.8) circle (0.24) (78.3,-13.7) circle (0.24) (80.1,0.8) circle (0.24) (87.0,-6.2) circle (0.24) (87.6,-6.6) circle (0.24) (74.4,-4.2) circle (0.24) (79.4,12.8) circle (0.24) (79.7,-4.0) circle (0.24) (78.8,4.4) circle (0.24) (98.5,8.5) circle (0.24)
    (89.1,0.5) circle (0.24) (100.2,10.3) circle (0.24) (90.6,12.1) circle (0.24) (85.9,-2.5) circle (0.24) (70.6,6.1) circle (0.24) (78.8,13.8) circle (0.24) (77.8,5.6) circle (0.24) (87.1,-2.8) circle (0.24) (95.0,8.5) circle (0.24) (77.0,3.7) circle (0.24) (97.0,-0.8) circle (0.24) (85.3,-1.1) circle (0.24)
    (79.2,14.4) circle (0.24) (87.2,3.8) circle (0.24) (70.6,-13.3) circle (0.24) (85.6,0.7) circle (0.24) (70.6,-5.2) circle (0.24) (95.0,-13.4) circle (0.24) (89.1,-4.6) circle (0.24) (99.0,-2.4) circle (0.24) (98.2,-14.0) circle (0.24) (98.6,14.0) circle (0.24) (104.2,7.2) circle (0.24) (78.4,-4.6) circle (0.24)
    (105.1,-1.9) circle (0.24) (87.9,0.6) circle (0.24) (97.3,-2.5) circle (0.24) (96.6,11.2) circle (0.24) (96.3,-2.6) circle (0.24) (90.3,-11.6) circle (0.24) (81.4,-13.0) circle (0.24) (90.3,-6.6) circle (0.24) (76.5,0.3) circle (0.24) (89.4,11.8) circle (0.24) (78.3,-10.6) circle (0.24) (70.7,-3.4) circle (0.24)
    (77.7,-6.9) circle (0.24) (79.0,-2.5) circle (0.24) (96.3,-14.3) circle (0.24) (86.5,9.2) circle (0.24) (88.0,-11.0) circle (0.24) (72.9,9.7) circle (0.24) (76.5,-8.6) circle (0.24) (99.6,12.6) circle (0.24) (79.6,5.0) circle (0.24) (78.7,-11.5) circle (0.24) (98.4,-10.5) circle (0.24) (77.8,-9.5) circle (0.24)
    (79.0,-12.2) circle (0.24) (78.2,0.9) circle (0.24) (71.4,1.9) circle (0.24) (91.5,7.8) circle (0.24) (78.9,0.9) circle (0.24) (89.3,4.0) circle (0.24) (97.1,2.0) circle (0.24) (85.9,8.1) circle (0.24) (78.2,-0.3) circle (0.24) (79.7,-0.1) circle (0.24) (87.0,-5.5) circle (0.24) (86.0,0.0) circle (0.24)
    (89.0,-10.6) circle (0.24) (96.5,8.5) circle (0.24) (70.7,-8.5) circle (0.24) (82.2,13.2) circle (0.24) (89.2,-12.5) circle (0.24) (96.6,4.1) circle (0.24) (80.0,13.8) circle (0.24) (98.4,-2.5) circle (0.24) (97.9,8.4) circle (0.24) (87.0,-0.6) circle (0.24) (79.5,-0.1) circle (0.24) (100.7,-0.6) circle (0.24)
    (86.9,-6.8) circle (0.24) (88.3,10.4) circle (0.24) (76.3,-10.5) circle (0.24) (88.4,-1.2) circle (0.24) (87.0,6.3) circle (0.24) (89.3,9.6) circle (0.24) (105.0,-3.2) circle (0.24) (86.2,7.0) circle (0.24) (88.2,11.2) circle (0.24) (86.3,-1.2) circle (0.24) (88.6,-9.0) circle (0.24) (78.5,-11.9) circle (0.24)
    (95.6,0.5) circle (0.24) (89.4,2.3) circle (0.24) (86.1,9.5) circle (0.24) (94.7,-11.4) circle (0.24) (97.2,14.3) circle (0.24) (91.5,-8.2) circle (0.24) (70.9,13.6) circle (0.24) (79.7,10.6) circle (0.24) (91.8,3.1) circle (0.24) (88.6,-8.0) circle (0.24) (98.4,-4.1) circle (0.24) (78.8,10.2) circle (0.24)
    (90.9,-0.8) circle (0.24) (95.5,0.2) circle (0.24) (99.2,3.7) circle (0.24) (80.6,0.3) circle (0.24) (89.1,7.6) circle (0.24) (88.0,-7.1) circle (0.24) (98.0,2.4) circle (0.24) (79.5,4.0) circle (0.24) (86.9,-6.3) circle (0.24) (87.6,-4.2) circle (0.24) (77.4,0.3) circle (0.24) (97.2,-13.5) circle (0.24)
    (72.0,11.7) circle (0.24) (89.6,6.0) circle (0.24) (98.9,-10.2) circle (0.24) (70.8,1.7) circle (0.24) (90.7,-4.8) circle (0.24) (97.4,10.6) circle (0.24) (87.7,5.9) circle (0.24) (87.7,4.7) circle (0.24) (95.3,-6.5) circle (0.24) (78.5,-0.5) circle (0.24) (81.2,-0.4) circle (0.24) (88.4,-0.2) circle (0.24)
    (96.5,9.6) circle (0.24) (89.5,9.8) circle (0.24) (79.4,-4.0) circle (0.24) (87.9,-8.0) circle (0.24) (88.5,-4.1) circle (0.24) (103.1,9.9) circle (0.24) (100.5,13.0) circle (0.24) (88.9,11.4) circle (0.24) (81.9,8.0) circle (0.24) (77.3,-12.6) circle (0.24) (95.6,9.3) circle (0.24) (98.7,4.9) circle (0.24)
    (76.7,-5.0) circle (0.24) (72.6,-5.0) circle (0.24) (80.2,-0.6) circle (0.24) (75.4,0.2) circle (0.24) (98.4,-12.7) circle (0.24) (79.6,7.3) circle (0.24) (85.7,-2.8) circle (0.24) (88.7,13.1) circle (0.24) (77.2,1.4) circle (0.24) (80.7,1.6) circle (0.24) (85.5,-11.2) circle (0.24) (75.6,-9.1) circle (0.24)
    (79.1,-10.2) circle (0.24) (85.9,-3.9) circle (0.24) (88.5,-8.4) circle (0.24) (88.3,13.3) circle (0.24) (99.8,-10.7) circle (0.24) (85.8,-10.2) circle (0.24) (89.7,-10.9) circle (0.24) (90.7,9.5) circle (0.24) (77.8,1.6) circle (0.24) (81.9,-0.4) circle (0.24) (88.8,-8.5) circle (0.24) (77.0,8.7) circle (0.24)
    (78.6,-13.5) circle (0.24) (87.9,-0.8) circle (0.24) (90.0,0.1) circle (0.24) (90.8,-6.2) circle (0.24) (77.5,6.8) circle (0.24) (71.1,-9.7) circle (0.24) (88.8,13.2) circle (0.24) (87.6,-13.2) circle (0.24) (78.1,-12.2) circle (0.24) (90.8,5.8) circle (0.24) (97.2,-6.9) circle (0.24) (86.4,8.5) circle (0.24)
    (101.1,-14.0) circle (0.24) (86.8,9.9) circle (0.24) (78.4,7.3) circle (0.24) (89.4,-2.5) circle (0.24) (94.7,5.5) circle (0.24) (89.5,-11.6) circle (0.24) (98.7,10.3) circle (0.24) (77.0,0.4) circle (0.24) (95.9,-5.7) circle (0.24) (85.2,-0.5) circle (0.24) (95.2,4.7) circle (0.24) (87.8,13.3) circle (0.24)
    (88.6,-6.6) circle (0.24) (88.4,-5.6) circle (0.24) (78.7,-11.2) circle (0.24) (88.7,-11.9) circle (0.24) (89.7,7.1) circle (0.24) (98.9,-2.4) circle (0.24) (80.4,10.7) circle (0.24) (97.3,3.4) circle (0.24) (89.1,-9.6) circle (0.24) (99.8,-5.4) circle (0.24) (77.5,7.5) circle (0.24) (70.7,-4.1) circle (0.24)
    (87.8,9.9) circle (0.24) (88.9,-5.2) circle (0.24) (78.2,3.3) circle (0.24) (88.3,-7.5) circle (0.24) (95.8,-10.4) circle (0.24) (89.6,-4.9) circle (0.24) (87.2,7.2) circle (0.24) (79.6,-6.9) circle (0.24) (77.7,11.4) circle (0.24) (96.2,1.6) circle (0.24) (88.7,4.9) circle (0.24) (88.0,-8.2) circle (0.24)
    (104.4,-0.9) circle (0.24) (99.3,8.6) circle (0.24) (94.3,-14.3) circle (0.24) (79.8,9.3) circle (0.24) (94.2,6.1) circle (0.24) (98.0,5.5) circle (0.24) (90.2,-7.5) circle (0.24) (87.1,-13.1) circle (0.24) (86.2,11.2) circle (0.24) (98.3,1.2) circle (0.24) (95.8,10.1) circle (0.24) (94.1,-0.6) circle (0.24)
    (79.9,-3.0) circle (0.24) (86.5,2.5) circle (0.24) (80.1,-11.2) circle (0.24) (81.1,-7.2) circle (0.24) (88.1,-7.1) circle (0.24) (87.1,-9.9) circle (0.24) (101.3,8.8) circle (0.24) (72.1,4.2) circle (0.24) (87.2,10.7) circle (0.24) (88.5,-0.2) circle (0.24) (96.9,-6.4) circle (0.24) (88.6,-14.3) circle (0.24)
    (88.7,6.3) circle (0.24) (77.8,2.4) circle (0.24) (88.3,11.0) circle (0.24) (88.5,5.2) circle (0.24) (87.7,14.3) circle (0.24) (96.2,5.8) circle (0.24) (98.6,-10.2) circle (0.24) (98.7,-4.3) circle (0.24) (86.7,11.4) circle (0.24) (76.7,-14.2) circle (0.24) (77.5,-9.3) circle (0.24) (87.2,4.4) circle (0.24)
    (84.9,5.4) circle (0.24) (82.0,9.7) circle (0.24) (80.8,5.7) circle (0.24) (96.7,-7.4) circle (0.24) (98.1,13.6) circle (0.24) (76.1,-3.4) circle (0.24) (98.7,8.3) circle (0.24) (82.3,-11.5) circle (0.24) (95.8,7.6) circle (0.24) (77.5,-8.3) circle (0.24) (79.5,6.3) circle (0.24) (94.4,-13.2) circle (0.24)
    (88.7,12.3) circle (0.24) (97.4,-13.5) circle (0.24) (89.7,-4.8) circle (0.24) (100.1,0.5) circle (0.24) (77.1,5.9) circle (0.24) (79.4,7.7) circle (0.24) (89.2,-2.3) circle (0.24) (98.3,11.3) circle (0.24) (77.8,-8.1) circle (0.24) (93.5,-2.4) circle (0.24) (71.3,8.9) circle (0.24) (88.9,-13.5) circle (0.24)
    (84.7,-0.6) circle (0.24) (79.5,0.8) circle (0.24) (88.2,9.0) circle (0.24) (96.3,6.1) circle (0.24) (77.5,-13.9) circle (0.24) (77.2,-4.2) circle (0.24) (96.3,10.8) circle (0.24) (75.7,0.7) circle (0.24) (74.9,-6.3) circle (0.24) (86.2,1.8) circle (0.24) (105.3,13.2) circle (0.24) (90.2,8.3) circle (0.24)
    (87.5,12.1) circle (0.24) (89.7,-1.2) circle (0.24) (78.5,3.9) circle (0.24) (102.9,6.2) circle (0.24) (85.3,-4.3) circle (0.24) (88.8,14.2) circle (0.24) (87.7,9.9) circle (0.24) (87.9,7.3) circle (0.24) (88.7,10.8) circle (0.24) (98.2,12.4) circle (0.24) (71.2,-5.5) circle (0.24) (100.2,5.0) circle (0.24)
    (98.7,-0.5) circle (0.24) (78.1,-7.4) circle (0.24) (96.2,-10.6) circle (0.24) (80.5,12.3) circle (0.24) (96.9,13.3) circle (0.24) (104.7,-7.1) circle (0.24) (97.6,-5.3) circle (0.24) (86.3,1.7) circle (0.24) (79.8,10.8) circle (0.24) (99.1,-1.4) circle (0.24) (85.6,8.0) circle (0.24) (95.4,-14.2) circle (0.24)
    (95.7,7.5) circle (0.24) (100.0,11.9) circle (0.24) (85.7,-3.5) circle (0.24) (89.0,-5.9) circle (0.24) (71.1,-13.9) circle (0.24) (89.9,-7.8) circle (0.24) (88.9,-13.4) circle (0.24) (88.3,9.8) circle (0.24) (96.9,0.7) circle (0.24) (97.9,1.1) circle (0.24) (96.3,5.2) circle (0.24) (88.9,9.7) circle (0.24)
    (80.0,-3.1) circle (0.24) (96.1,-13.6) circle (0.24) (71.5,-4.7) circle (0.24) (87.5,1.9) circle (0.24) (96.8,-12.0) circle (0.24) (78.5,-14.2) circle (0.24) (99.4,-1.3) circle (0.24) (87.9,-7.8) circle (0.24) (104.9,-7.9) circle (0.24) (71.4,-5.5) circle (0.24) (75.4,11.6) circle (0.24) (98.3,-14.3) circle (0.24)
    (89.1,-14.2) circle (0.24) (88.9,-14.4) circle (0.24) (105.3,12.4) circle (0.24) (97.9,5.0) circle (0.24) (87.6,9.6) circle (0.24) (97.1,1.8) circle (0.24) (98.6,-0.7) circle (0.24) (87.5,6.8) circle (0.24) (85.2,10.2) circle (0.24) (78.1,0.4) circle (0.24) (97.5,2.8) circle (0.24) (87.8,14.3) circle (0.24)
    (80.3,11.4) circle (0.24) (88.9,9.5) circle (0.24) (98.2,4.1) circle (0.24) (88.0,-7.2) circle (0.24) (86.1,-9.1) circle (0.24) (70.6,2.2) circle (0.24) (77.0,-5.0) circle (0.24) (96.7,-13.5) circle (0.24) (91.4,3.6) circle (0.24) (99.6,9.9) circle (0.24) (81.4,-13.3) circle (0.24) (76.5,-3.6) circle (0.24)
    (95.8,-13.0) circle (0.24) (93.9,12.3) circle (0.24) (94.7,-2.6) circle (0.24) (72.4,-4.3) circle (0.24) (96.9,2.4) circle (0.24) (88.4,-0.1) circle (0.24) (81.0,5.2) circle (0.24) (85.7,-8.4) circle (0.24) (77.1,-6.2) circle (0.24) (78.6,-4.9) circle (0.24) (98.2,-1.2) circle (0.24) (77.5,13.0) circle (0.24)
    (88.5,6.4) circle (0.24) (96.5,-7.2) circle (0.24) (91.4,12.6) circle (0.24) (95.7,-4.0) circle (0.24) (97.6,-11.5) circle (0.24) (98.2,13.9) circle (0.24) (89.3,-1.4) circle (0.24) (99.1,0.1) circle (0.24) (89.5,-13.9) circle (0.24) (98.9,-9.8) circle (0.24) (96.4,3.5) circle (0.24) (71.5,-9.4) circle (0.24)
    (89.4,-1.0) circle (0.24) (79.4,-12.8) circle (0.24) (97.0,11.3) circle (0.24) (88.4,-3.5) circle (0.24) (79.4,-8.5) circle (0.24) (85.7,-3.9) circle (0.24) (80.0,11.5) circle (0.24) (96.0,-13.7) circle (0.24) (77.3,-12.5) circle (0.24) (89.5,10.0) circle (0.24) (96.2,-9.1) circle (0.24) (81.9,4.8) circle (0.24)
    (93.8,-9.6) circle (0.24) (87.5,12.7) circle (0.24) (95.8,12.8) circle (0.24) (89.6,-8.0) circle (0.24) (86.1,14.0) circle (0.24) (72.1,4.6) circle (0.24) (94.5,11.8) circle (0.24) (87.1,12.6) circle (0.24) (94.9,-3.7) circle (0.24) (95.1,-9.6) circle (0.24) (86.5,-0.7) circle (0.24) (87.2,3.4) circle (0.24)
    (79.0,10.9) circle (0.24) (88.5,-14.3) circle (0.24) (86.4,-1.4) circle (0.24) (72.3,9.5) circle (0.24) (97.1,-5.3) circle (0.24) (88.7,0.2) circle (0.24) (90.4,-10.5) circle (0.24) (94.0,11.5) circle (0.24) (72.4,14.3) circle (0.24) (91.4,-4.7) circle (0.24) (97.9,-9.2) circle (0.24) (76.6,3.3) circle (0.24)
    (105.3,5.5) circle (0.24) (86.7,9.8) circle (0.24) (95.5,-10.0) circle (0.24) (97.3,1.8) circle (0.24) (79.1,8.8) circle (0.24) (78.1,-10.2) circle (0.24) (76.5,-13.0) circle (0.24) (89.6,-12.0) circle (0.24) (82.0,-9.8) circle (0.24) (100.2,2.8) circle (0.24) (97.2,-4.9) circle (0.24) (88.8,-11.5) circle (0.24)
    (95.5,6.2) circle (0.24) (91.5,-14.1) circle (0.24) (96.8,-2.8) circle (0.24) (80.5,0.1) circle (0.24) (78.8,7.1) circle (0.24) (88.2,-4.3) circle (0.24) (90.5,-3.2) circle (0.24) (90.1,-9.1) circle (0.24) (101.1,11.7) circle (0.24) (89.6,3.6) circle (0.24) (96.9,-5.1) circle (0.24) (98.1,4.0) circle (0.24)
    (80.9,-6.6) circle (0.24) (79.5,-0.3) circle (0.24) (86.0,-7.9) circle (0.24) (96.7,11.0) circle (0.24) (98.7,13.9) circle (0.24) (89.2,10.7) circle (0.24) (80.0,-1.0) circle (0.24) (80.6,4.8) circle (0.24) (104.6,-8.0) circle (0.24) (80.4,-0.3) circle (0.24) (79.5,-1.6) circle (0.24) (87.3,-4.1) circle (0.24)
    (99.1,-3.1) circle (0.24) (89.0,1.5) circle (0.24) (105.4,-1.0) circle (0.24) (90.3,1.5) circle (0.24) (97.2,-5.2) circle (0.24) (88.0,11.2) circle (0.24) (86.8,-13.8) circle (0.24) (84.9,-14.0) circle (0.24) (86.0,-11.4) circle (0.24) (87.1,-8.0) circle (0.24) (105.0,-5.9) circle (0.24) (78.0,8.9) circle (0.24)
    (87.0,5.0) circle (0.24) (97.3,0.2) circle (0.24) (82.1,5.9) circle (0.24) (88.9,-11.2) circle (0.24) (95.7,5.8) circle (0.24) (86.8,-10.0) circle (0.24) (78.9,-8.6) circle (0.24) (94.5,-13.0) circle (0.24) (90.9,6.2) circle (0.24) (80.7,5.9) circle (0.24) (88.4,12.1) circle (0.24) (85.8,-0.2) circle (0.24)
    (87.7,8.2) circle (0.24) (77.9,9.5) circle (0.24) (78.7,-6.5) circle (0.24) (76.1,-1.7) circle (0.24) (96.4,13.3) circle (0.24) (79.6,-12.5) circle (0.24) (81.6,1.1) circle (0.24) (78.5,-2.4) circle (0.24) (80.9,13.8) circle (0.24) (95.9,5.7) circle (0.24) (88.9,5.9) circle (0.24) (95.6,10.6) circle (0.24)
    (88.3,10.3) circle (0.24) (86.7,14.1) circle (0.24) (96.7,7.0) circle (0.24) (87.8,-2.3) circle (0.24) (97.6,6.1) circle (0.24) (97.6,9.8) circle (0.24) (88.8,-5.5) circle (0.24) (76.7,-7.7) circle (0.24) (97.2,-3.8) circle (0.24) (91.1,3.1) circle (0.24) (78.0,3.3) circle (0.24) (97.1,-8.1) circle (0.24)
    (105.3,-8.9) circle (0.24) (79.0,-10.2) circle (0.24) (78.8,-13.9) circle (0.24) (95.3,-0.6) circle (0.24) (78.0,10.5) circle (0.24) (89.4,-2.5) circle (0.24) (78.6,-10.7) circle (0.24) (78.2,-2.4) circle (0.24) (97.1,-1.3) circle (0.24) (77.7,5.7) circle (0.24) (85.7,10.8) circle (0.24) (75.9,-9.9) circle (0.24)
    (81.4,-13.8) circle (0.24) (89.2,-8.3) circle (0.24) (79.6,-9.4) circle (0.24) (85.2,2.0) circle (0.24) (80.8,6.0) circle (0.24) (90.1,-7.3) circle (0.24) (78.2,2.3) circle (0.24) (89.5,-11.0) circle (0.24) (85.0,1.3) circle (0.24) (97.5,10.8) circle (0.24) (85.7,-0.6) circle (0.24) (89.4,-7.7) circle (0.24)
    (90.3,0.6) circle (0.24) (97.5,13.1) circle (0.24) (77.6,1.1) circle (0.24) (90.4,11.9) circle (0.24) (79.4,-8.9) circle (0.24) (89.8,-14.4) circle (0.24) (100.9,-8.0) circle (0.24) (85.7,-9.6) circle (0.24) (98.6,-5.7) circle (0.24) (70.6,9.5) circle (0.24) (85.9,11.8) circle (0.24) (80.5,-9.1) circle (0.24)
    (79.7,2.0) circle (0.24) (81.5,10.8) circle (0.24) (76.4,-1.4) circle (0.24) (82.3,-3.2) circle (0.24) (88.6,9.8) circle (0.24) (90.3,4.3) circle (0.24) (96.1,5.0) circle (0.24) (76.5,-5.8) circle (0.24) (89.7,1.4) circle (0.24) (99.0,0.4) circle (0.24) (78.0,-6.2) circle (0.24) (90.6,7.4) circle (0.24)
    (86.3,-2.0) circle (0.24) (98.3,-10.9) circle (0.24) (89.7,7.1) circle (0.24) (103.1,-2.6) circle (0.24) (81.6,-8.2) circle (0.24) (70.7,3.5) circle (0.24) (78.1,-12.7) circle (0.24) (88.1,-14.3) circle (0.24) (80.6,3.2) circle (0.24) (96.3,-1.9) circle (0.24) (89.0,4.8) circle (0.24) (87.4,4.7) circle (0.24)
    (89.9,-9.3) circle (0.24) (85.6,-0.6) circle (0.24) (90.6,4.6) circle (0.24) (95.3,-7.2) circle (0.24) (72.1,0.2) circle (0.24) (86.1,-10.6) circle (0.24) (84.3,-6.4) circle (0.24) (84.1,-8.5) circle (0.24) (89.1,-13.5) circle (0.24) (78.1,10.5) circle (0.24) (87.0,-10.2) circle (0.24) (88.3,-4.3) circle (0.24)
    (87.6,11.7) circle (0.24) (98.4,-4.1) circle (0.24) (94.7,-7.0) circle (0.24) (77.3,11.8) circle (0.24) (88.1,5.4) circle (0.24) (86.0,-10.7) circle (0.24) (79.0,-4.1) circle (0.24) (97.8,-0.4) circle (0.24) (78.5,2.6) circle (0.24) (70.9,7.0) circle (0.24) (96.3,2.5) circle (0.24) (96.4,8.7) circle (0.24)
    (80.5,8.6) circle (0.24) (96.7,-6.4) circle (0.24) (85.5,-12.9) circle (0.24) (87.8,-6.4) circle (0.24) (96.5,4.6) circle (0.24) (86.7,-8.2) circle (0.24) (87.0,8.7) circle (0.24) (77.9,-5.4) circle (0.24) (86.6,-7.9) circle (0.24) (87.2,2.9) circle (0.24) (97.9,-3.4) circle (0.24) (89.4,-4.9) circle (0.24)
    (87.5,1.8) circle (0.24) (80.8,0.7) circle (0.24) (88.0,-2.3) circle (0.24) (104.9,-14.2) circle (0.24) (80.3,-1.2) circle (0.24) (79.7,8.9) circle (0.24) (72.4,10.6) circle (0.24) (88.5,13.6) circle (0.24) (98.0,5.9) circle (0.24) (81.5,-5.7) circle (0.24) (80.1,10.8) circle (0.24) (98.4,-0.0) circle (0.24)
    (85.9,3.9) circle (0.24) (98.1,-7.3) circle (0.24) (97.5,-13.9) circle (0.24) (80.2,-1.4) circle (0.24) (71.5,-3.7) circle (0.24) (77.0,-11.9) circle (0.24) (80.5,-4.8) circle (0.24) (90.8,3.6) circle (0.24) (77.8,-9.0) circle (0.24) (86.1,4.1) circle (0.24) (89.7,-6.1) circle (0.24) (88.3,-12.8) circle (0.24)
    (96.7,2.8) circle (0.24) (79.8,-0.1) circle (0.24) (70.7,11.7) circle (0.24) (93.9,12.4) circle (0.24) (91.6,0.2) circle (0.24) (96.7,6.7) circle (0.24) (94.3,-2.6) circle (0.24) (88.3,7.7) circle (0.24) (79.8,13.0) circle (0.24) (101.0,-3.5) circle (0.24) (78.2,-3.0) circle (0.24) (95.7,-3.6) circle (0.24)
    (96.2,6.4) circle (0.24) (78.4,-12.3) circle (0.24) (71.4,11.9) circle (0.24) (99.2,8.7) circle (0.24) (90.8,5.4) circle (0.24) (85.0,10.1) circle (0.24) (95.6,-2.1) circle (0.24) (79.5,13.0) circle (0.24) (105.2,4.1) circle (0.24) (94.6,11.5) circle (0.24) (88.5,-9.1) circle (0.24) (87.6,1.6) circle (0.24)
    (96.3,11.0) circle (0.24) (99.0,-3.7) circle (0.24) (87.6,7.6) circle (0.24) (85.0,-3.3) circle (0.24) (94.1,-2.2) circle (0.24) (105.3,-4.5) circle (0.24) (75.7,-6.5) circle (0.24) (76.6,0.7) circle (0.24) (88.2,0.0) circle (0.24) (76.7,-10.9) circle (0.24) (86.2,-4.8) circle (0.24) (76.9,4.5) circle (0.24)
    (96.2,-12.8) circle (0.24) (89.8,7.7) circle (0.24) (98.7,-6.8) circle (0.24) (87.8,-2.9) circle (0.24) (96.6,-13.2) circle (0.24) (77.9,9.8) circle (0.24) (96.9,-8.6) circle (0.24) (84.9,-2.4) circle (0.24) (104.6,11.5) circle (0.24) (96.9,10.0) circle (0.24) (80.5,2.2) circle (0.24) (104.5,10.0) circle (0.24)
    (77.2,13.7) circle (0.24) (103.6,-5.7) circle (0.24) (77.4,13.6) circle (0.24) (96.5,7.0) circle (0.24) (78.9,-9.4) circle (0.24) (96.7,14.0) circle (0.24) (96.3,3.0) circle (0.24) (77.0,-7.0) circle (0.24) (80.3,-13.4) circle (0.24) (81.3,11.5) circle (0.24) (71.5,8.4) circle (0.24) (71.0,-12.4) circle (0.24)
    (90.9,8.2) circle (0.24) (77.2,8.5) circle (0.24) (79.5,-6.2) circle (0.24) (99.5,-9.4) circle (0.24) (98.8,-0.1) circle (0.24) (70.7,2.9) circle (0.24) (77.9,5.3) circle (0.24) (78.3,-2.9) circle (0.24) (72.2,4.4) circle (0.24) (76.2,-11.9) circle (0.24) (99.5,11.4) circle (0.24) (99.7,2.6) circle (0.24)
    (77.4,-8.1) circle (0.24) (77.6,-13.7) circle (0.24) (100.7,-13.7) circle (0.24) (89.1,-10.1) circle (0.24) (78.3,-6.7) circle (0.24) (87.6,-6.8) circle (0.24) (98.9,1.1) circle (0.24) (98.9,-4.4) circle (0.24) (105.4,-0.5) circle (0.24) (88.5,-4.2) circle (0.24) (79.0,13.7) circle (0.24) (86.4,-7.9) circle (0.24)
    (87.4,9.4) circle (0.24) (85.4,-0.6) circle (0.24) (81.2,7.8) circle (0.24) (81.1,8.1) circle (0.24) (97.6,3.2) circle (0.24) (76.5,10.0) circle (0.24) (71.1,-2.9) circle (0.24) (80.0,-10.9) circle (0.24) (82.0,13.3) circle (0.24) (79.1,-8.5) circle (0.24) (80.0,1.0) circle (0.24) (80.5,6.2) circle (0.24)
    (78.4,0.4) circle (0.24) (89.1,-6.9) circle (0.24) (80.0,-5.7) circle (0.24) (95.8,6.1) circle (0.24) (85.7,12.9) circle (0.24) (77.6,-9.5) circle (0.24) (88.3,-0.5) circle (0.24) (87.8,-12.7) circle (0.24) (82.1,7.5) circle (0.24) (89.4,-10.7) circle (0.24) (89.9,6.1) circle (0.24) (99.6,-2.4) circle (0.24)
    (85.9,-0.7) circle (0.24) (75.7,-12.9) circle (0.24) (70.5,11.6) circle (0.24) (98.2,-5.8) circle (0.24) (77.8,6.1) circle (0.24) (78.8,-8.4) circle (0.24) (78.4,-7.6) circle (0.24) (79.8,-9.5) circle (0.24) (72.3,-14.3) circle (0.24) (94.1,-0.6) circle (0.24) (85.5,-12.9) circle (0.24) (76.0,-8.4) circle (0.24)
    (82.1,3.3) circle (0.24) (77.5,5.6) circle (0.24) (87.9,8.1) circle (0.24) (86.2,-13.2) circle (0.24) (80.7,-9.2) circle (0.24) (79.2,11.1) circle (0.24) (78.8,12.6) circle (0.24) (86.8,13.2) circle (0.24) (95.2,-13.0) circle (0.24) (76.8,14.3) circle (0.24) (87.4,8.8) circle (0.24) (78.6,12.5) circle (0.24)
    (77.4,12.0) circle (0.24) (98.8,5.4) circle (0.24) (88.7,11.0) circle (0.24) (86.8,8.7) circle (0.24) (95.4,-11.3) circle (0.24) (97.6,-12.1) circle (0.24) (85.6,-4.4) circle (0.24) (78.6,-3.4) circle (0.24) (95.5,-7.8) circle (0.24) (75.9,-11.5) circle (0.24) (77.1,6.3) circle (0.24) (76.5,11.4) circle (0.24)
    (97.2,4.4) circle (0.24) (72.2,8.2) circle (0.24) (77.0,4.4) circle (0.24) (86.9,5.4) circle (0.24) (72.8,-12.6) circle (0.24) (75.9,-8.5) circle (0.24) (98.3,-12.1) circle (0.24) (85.8,-11.8) circle (0.24) (87.5,10.0) circle (0.24) (70.6,-12.3) circle (0.24) (89.8,-4.9) circle (0.24) (77.6,13.5) circle (0.24)
    (76.3,0.0) circle (0.24) (77.7,14.1) circle (0.24) (77.6,2.5) circle (0.24) (98.1,-10.6) circle (0.24) (70.7,6.4) circle (0.24) (98.5,-1.7) circle (0.24) (84.9,-7.4) circle (0.24) (87.7,-2.4) circle (0.24) (86.4,13.0) circle (0.24) (86.6,-10.0) circle (0.24) (88.0,-5.7) circle (0.24) (89.1,-14.0) circle (0.24)
    (87.9,11.4) circle (0.24) (87.5,-13.0) circle (0.24) (71.0,13.1) circle (0.24) (78.2,-10.9) circle (0.24) (86.7,2.6) circle (0.24) (103.8,-9.6) circle (0.24) (100.2,12.8) circle (0.24) (86.5,4.2) circle (0.24) (94.4,4.0) circle (0.24) (104.2,7.1) circle (0.24) (93.2,-1.2) circle (0.24) (89.5,12.5) circle (0.24)
    (79.7,-5.0) circle (0.24) (71.7,-0.1) circle (0.24) (88.8,-5.4) circle (0.24) (77.7,-3.9) circle (0.24) (78.9,7.1) circle (0.24) (77.7,13.6) circle (0.24) (80.4,-8.4) circle (0.24) (95.4,9.8) circle (0.24) (80.8,0.4) circle (0.24) (79.8,1.7) circle (0.24) (88.3,-9.5) circle (0.24) (86.8,-3.6) circle (0.24)
    (77.0,1.2) circle (0.24) (77.3,0.8) circle (0.24) (82.6,0.5) circle (0.24) (87.1,-4.3) circle (0.24) (96.2,-12.6) circle (0.24) (78.5,6.3) circle (0.24) (78.5,-0.6) circle (0.24) (86.8,5.9) circle (0.24) (85.2,-11.0) circle (0.24) (97.4,-8.5) circle (0.24) (97.7,1.3) circle (0.24) (94.6,-11.3) circle (0.24)
    (90.0,13.5) circle (0.24) (90.7,14.2) circle (0.24) (81.3,-0.2) circle (0.24) (89.0,5.0) circle (0.24) (99.5,-7.6) circle (0.24) (105.2,11.5) circle (0.24) (98.3,11.2) circle (0.24) (97.3,-6.3) circle (0.24) (76.1,1.2) circle (0.24) (88.9,0.9) circle (0.24) (88.8,11.0) circle (0.24) (94.7,-10.4) circle (0.24)
    (87.2,1.7) circle (0.24) (78.1,-11.8) circle (0.24) (99.6,10.1) circle (0.24) (78.8,0.8) circle (0.24) (78.3,6.5) circle (0.24) (86.9,2.7) circle (0.24) (86.1,-5.1) circle (0.24) (87.4,-1.8) circle (0.24) (87.5,-5.8) circle (0.24) (89.2,2.0) circle (0.24) (88.7,-1.3) circle (0.24) (95.2,1.6) circle (0.24)
    (96.2,3.0) circle (0.24) (85.9,-6.0) circle (0.24) (77.3,0.9) circle (0.24) (84.5,-3.9) circle (0.24) (88.7,12.3) circle (0.24) (87.6,5.5) circle (0.24) (89.0,-3.9) circle (0.24) (78.4,-10.5) circle (0.24) (85.9,-14.2) circle (0.24) (90.7,11.0) circle (0.24) (90.6,-6.8) circle (0.24) (78.7,-5.8) circle (0.24)
    (88.3,5.6) circle (0.24) (81.4,-10.1) circle (0.24) (77.6,-5.4) circle (0.24) (76.6,7.0) circle (0.24) (95.4,-8.5) circle (0.24) (79.3,-11.8) circle (0.24) (77.2,9.9) circle (0.24) (104.6,9.1) circle (0.24) (77.5,2.8) circle (0.24) (100.0,-10.2) circle (0.24) (88.5,-13.3) circle (0.24) (95.2,-3.5) circle (0.24)
    (76.3,1.0) circle (0.24) (88.4,-4.8) circle (0.24) (86.9,-2.3) circle (0.24) (88.0,6.7) circle (0.24) (80.9,-0.5) circle (0.24) (97.6,0.6) circle (0.24) (98.4,14.0) circle (0.24) (89.8,-10.8) circle (0.24) (105.0,1.5) circle (0.24) (87.5,-11.0) circle (0.24) (87.9,-12.6) circle (0.24) (105.3,4.5) circle (0.24)
    (77.6,3.9) circle (0.24) (77.7,-2.2) circle (0.24) (76.7,-12.6) circle (0.24) (89.5,0.4) circle (0.24) (94.7,-3.5) circle (0.24) (77.2,6.1) circle (0.24) (85.5,-1.1) circle (0.24) (103.3,5.7) circle (0.24) (87.1,10.1) circle (0.24) (80.8,-9.6) circle (0.24) (100.2,12.0) circle (0.24) (89.2,12.2) circle (0.24)
    (95.7,-8.1) circle (0.24) (96.7,-9.6) circle (0.24) (97.2,5.6) circle (0.24) (75.8,-1.7) circle (0.24) (79.6,9.6) circle (0.24) (96.5,2.3) circle (0.24) (89.1,9.9) circle (0.24) (85.4,-0.7) circle (0.24) (81.1,8.1) circle (0.24) (91.3,0.3) circle (0.24) (105.3,-2.2) circle (0.24) (93.4,3.3) circle (0.24)
    (78.0,-12.2) circle (0.24) (77.6,-2.2) circle (0.24) (89.2,1.1) circle (0.24) (86.9,12.4) circle (0.24) (88.6,7.2) circle (0.24) (88.1,-6.7) circle (0.24) (90.9,-10.8) circle (0.24) (79.1,10.4) circle (0.24) (80.1,1.4) circle (0.24) (71.1,-8.6) circle (0.24) (105.0,8.6) circle (0.24) (96.3,7.8) circle (0.24)
    (95.5,12.7) circle (0.24) (80.7,0.7) circle (0.24) (89.9,-12.5) circle (0.24) (77.1,6.3) circle (0.24) (96.9,-6.1) circle (0.24) (79.2,4.6) circle (0.24) (75.5,-13.9) circle (0.24) (75.8,-2.7) circle (0.24) (81.8,9.1) circle (0.24) (105.1,-1.4) circle (0.24) (89.7,4.8) circle (0.24) (97.8,-3.4) circle (0.24)
    (86.4,-10.0) circle (0.24) (77.7,-2.5) circle (0.24) (95.9,-11.3) circle (0.24) (89.9,-8.4) circle (0.24) (88.0,-0.9) circle (0.24) (96.5,-0.3) circle (0.24) (77.1,-9.8) circle (0.24) (81.2,-3.8) circle (0.24) (97.9,11.3) circle (0.24) (86.2,-7.9) circle (0.24) (98.4,-9.2) circle (0.24) (86.8,-2.3) circle (0.24)
    (91.2,-2.5) circle (0.24) (100.3,4.8) circle (0.24) (81.0,8.7) circle (0.24) (87.5,-10.7) circle (0.24) (84.7,10.3) circle (0.24) (98.6,6.6) circle (0.24) (77.2,5.2) circle (0.24) (105.3,4.8) circle (0.24) (88.1,-9.7) circle (0.24) (77.2,3.0) circle (0.24) (89.8,-3.4) circle (0.24) (78.8,10.8) circle (0.24)
    (93.3,0.3) circle (0.24) (99.1,-3.0) circle (0.24) (89.5,10.1) circle (0.24) (95.2,3.3) circle (0.24) (81.0,-0.4) circle (0.24) (77.4,-4.2) circle (0.24) (86.3,8.8) circle (0.24) (80.0,-11.1) circle (0.24) (87.8,-13.9) circle (0.24) (77.3,9.7) circle (0.24) (87.6,2.9) circle (0.24) (72.8,-9.1) circle (0.24)
    (81.0,3.9) circle (0.24) (97.7,9.6) circle (0.24) (98.0,6.0) circle (0.24) (70.8,-7.0) circle (0.24) (79.4,7.3) circle (0.24) (76.4,2.3) circle (0.24) (80.5,5.5) circle (0.24) (85.6,0.7) circle (0.24) (96.8,-12.5) circle (0.24) (78.6,-5.0) circle (0.24) (78.7,-7.1) circle (0.24) (88.0,7.8) circle (0.24)
    (87.0,-7.6) circle (0.24) (75.4,-10.5) circle (0.24) (81.7,-5.9) circle (0.24) (88.3,-13.9) circle (0.24) (95.9,-13.5) circle (0.24) (77.9,5.9) circle (0.24) (95.5,-13.4) circle (0.24) (90.6,-4.5) circle (0.24) (95.7,13.1) circle (0.24) (87.2,-12.8) circle (0.24) (81.0,-2.0) circle (0.24) (88.4,5.5) circle (0.24)
    (91.0,-4.9) circle (0.24) (89.7,-8.6) circle (0.24) (80.9,-3.6) circle (0.24) (98.9,-6.4) circle (0.24) (96.2,14.2) circle (0.24) (87.3,-2.8) circle (0.24) (88.8,-3.8) circle (0.24) (76.0,4.1) circle (0.24) (87.1,-13.8) circle (0.24) (96.9,0.5) circle (0.24) (84.4,-12.1) circle (0.24) (77.2,1.7) circle (0.24)
    (77.7,-10.1) circle (0.24) (81.9,-8.9) circle (0.24) (89.0,-11.3) circle (0.24) (101.0,-2.5) circle (0.24) (80.8,-5.9) circle (0.24) (79.6,11.5) circle (0.24) (87.3,-7.4) circle (0.24) (78.9,9.4) circle (0.24) (98.6,2.3) circle (0.24) (88.5,8.5) circle (0.24) (88.8,6.1) circle (0.24) (87.1,-5.8) circle (0.24)
    (75.8,13.6) circle (0.24) (91.2,-9.0) circle (0.24) (99.2,8.5) circle (0.24) (94.9,-9.9) circle (0.24) (80.2,6.3) circle (0.24) (95.5,-4.0) circle (0.24) (78.2,-6.2) circle (0.24) (96.2,3.7) circle (0.24) (97.3,-0.7) circle (0.24) (98.3,5.8) circle (0.24) (80.5,4.8) circle (0.24) (79.5,-1.1) circle (0.24)
    (87.9,5.7) circle (0.24) (85.7,2.9) circle (0.24) (78.3,-13.4) circle (0.24) (97.2,9.4) circle (0.24) (90.3,-13.4) circle (0.24) (91.3,0.6) circle (0.24) (97.6,-1.8) circle (0.24) (89.4,-8.6) circle (0.24) (75.8,12.6) circle (0.24) (90.6,-13.5) circle (0.24) (79.3,-13.8) circle (0.24) (75.4,-0.8) circle (0.24)
    (85.0,-3.6) circle (0.24) (78.4,-5.2) circle (0.24) (72.1,9.3) circle (0.24) (94.4,-1.2) circle (0.24) (87.7,14.2) circle (0.24) (76.7,-12.0) circle (0.24) (105.4,-0.9) circle (0.24) (97.6,13.8) circle (0.24) (88.5,7.8) circle (0.24) (78.5,9.0) circle (0.24) (98.5,10.9) circle (0.24) (78.8,-7.4) circle (0.24)
    (89.2,1.3) circle (0.24) (80.4,5.7) circle (0.24) (81.0,-14.3) circle (0.24) (89.5,-6.1) circle (0.24) (89.6,-13.6) circle (0.24) (87.1,5.0) circle (0.24) (85.4,10.0) circle (0.24) (85.3,2.1) circle (0.24) (80.3,14.1) circle (0.24) (81.1,2.3) circle (0.24) (86.2,-1.1) circle (0.24) (71.0,6.1) circle (0.24)
    (71.7,-1.9) circle (0.24) (94.9,12.3) circle (0.24) (95.0,4.2) circle (0.24) (81.5,9.0) circle (0.24) (97.3,0.4) circle (0.24) (70.8,-7.9) circle (0.24) (96.6,10.4) circle (0.24) (82.0,-1.3) circle (0.24) (88.8,-8.3) circle (0.24) (76.8,-1.6) circle (0.24) (87.4,-10.3) circle (0.24) (70.9,8.5) circle (0.24)
    (95.1,-4.9) circle (0.24) (78.4,-1.5) circle (0.24) (88.2,-12.1) circle (0.24) (86.8,-2.7) circle (0.24) (87.0,5.7) circle (0.24) (95.3,-9.5) circle (0.24) (101.0,-8.5) circle (0.24) (77.2,3.5) circle (0.24);
  \foreach \X/\N in {0/20, 44/200, 88/\num{2000}}{
    \draw[bkInk, line width=0.5pt] (\X-18,-15) rectangle (\X+18,15);
    \node[txd lab, anchor=north] at (\X,-16) {\N 個};}
\end{tikzpicture}
\bkfigcaption{電子を1個ずつスクリーンに落とす（⑤のプログラムで、届く確率を $|\psi|^2$ に比例させて作った図）。20個では、ばらばらの点にしか見えない。200個では、点の多い所と少ない所が見えてきて、2000個では、はっきりしたしまになる。}\label{ph:原-3:fig-ichikozutsu}
\end{bkfigure}

結果は、図\ref{ph:原-3:fig-ichikozutsu}と同じでした。
\begin{itemize}[leftmargin=1.2\zw, itemsep=2pt, topsep=2pt, label={\color{bkTheme}●}]
\item 電子は、1個ずつ、スクリーンの\textbf{1つの点}に届く（つぶとしてのふるまい）。
\item 1個がどこに届くかは、前もって決まっていない。ばらばらに見える。
\item たくさん集まると、\textbf{干渉のしま}ができる（波としてのふるまい）。しまの間かくは、波-13 ②の $\Delta x=L\lambda/d$ に、ド・ブロイ波長を入れた値になる（例題6）。
\end{itemize}
1個ずつ送っているので、電子どうしがぶつかって、しまを作ったのではありません。\textbf{1個の電子の波が、2つの道を同時に通り、自分自身と干渉した}と考えるしかありません。原-2 ⑧で見た光子1個ずつの干渉と、まったく同じです。光も電子も、「波か、つぶか」ではなく、その\textbf{両方}の性質をもっているのです。

\begin{phyomi}[確率の波]
電子の波の高さを $\psi$（プサイ）と書くと、電子がその場所で見つかる確率は、波の高さの2乗 $|\psi|^2$ に比例します。光の明るさが、電場の波の高さ（振幅）の2乗に比例するのと同じです。$|\psi|^2$ の大きい所（しまの明るい所）には電子が多く届き、0 の所（暗い所）には1個も届きません。$\psi$ をどう計算するかは、原-5 で学びます。
\end{phyomi}
\bkidx{かくりつのなみ}{確率の波}

\Needspace{8\baselineskip}\noindent\textbf{乱数で、電子を1個ずつ落とす}\par\nobreak
コンピュータで、$|\psi|^2$ に比例する確率で、電子を1個ずつ落としてみましょう（3-17）。しまの間かくは例題6 の \qty{5.4}{\micro m} とし、全体はゆるやかに暗くなるとします。$-10$〜\qty{10}{\micro m} の中で、届く位置の候補 $x$ を乱数で選び、0〜1 の乱数が「届きやすさ $P(x)$」より小さければ、そこに届いたとします（はねのけ法）。
\begin{phcode}[Python]
import random, math
random.seed(3)
s = 5.4          \# しまの間かく（\si{\micro m}）：例題6
w = 12.0         \# 全体の広がり（\si{\micro m}）
def P(x):        \# 届きやすさ（いちばん大きい所で1）
    return math.cos(math.pi * x / s)**2 * math.exp(-(x / w)**2)
def one():       \# 電子1個がスクリーンに届く位置
    while True:
        x = random.uniform(-10, 10)
        if random.random() < P(x):
            return x
for N in [10, 100, 10000]:
    hist = [0] * 10     \# 2 \si{\micro m} ずつ10の区間に分ける
    for i in range(N):
        hist[int((one() + 10) // 2)] += 1
    print(N, hist)
for k in range(10):     \# 10000個のときの棒グラフ
    print('\%3d' \% (-10 + 2 * k), '*' * (hist[k] // 50))
\end{phcode}
\begin{phcodeout}[実行結果]
10 [0, 0, 3, 1, 1, 2, 1, 2, 0, 0]
100 [3, 12, 16, 5, 13, 15, 3, 25, 8, 0]
10000 [418, 761, 1816, 305, 1720, 1612, 322, 1895, 742, 409]
-10 ********
 -8 ***************
 -6 ************************************
 -4 ******
 -2 **********************************
  0 ********************************
  2 ******
  4 *************************************
  6 **************
  8 ********
\end{phcodeout}
10個や100個では、どこが明るいのかよくわかりません。1万個になると、$-4$〜$-2$ と 2〜4 の区間（$\pm\qty{2.7}{\micro m}$ の暗い線をふくむ）がへこみ、しまがはっきり見えます。1個ずつの行き先は乱数でばらばらなのに、全体は決まった形になる。これが「確率の波」の意味です。乱数の種（\texttt{seed}）を変えると、1個ずつの行き先は変わりますが、1万個の形はほとんど変わりません。

\Needspace{14\baselineskip}\bksubsection{6}{はてな？の答え：電子顕微鏡}
\label{ph:原-3:sec6}
\label{ph:原-3:hatena}\bkhatenaref の答えです。波で物を見るとき、見分けられる細かさには限りがあります。レンズの口を通るときの回折で、1つの点の像も少し広がった円になるからです（波-13 ⑤）。見分けられる2点の間かくの目安は、
\pheq{\delta\approx\frac{0.61\,\lambda}{\alpha}}{m ＝ m ÷ rad}
です（$\alpha$ はレンズが試料を見こむ角の半分。小さいとき）。光の顕微鏡では、$\alpha$ を大きくしても（油にひたすレンズで見かけを大きくしても）、$\delta$ は波長の半分ほど、約 \qty{0.2}{\micro m} が限界です。

電子顕微鏡は、光のかわりに、電圧で加速した電子の波を使います。100 kV の電子の波長は \qty{3.7}{pm} で、緑の光（\qty{550}{nm}）の約15万分の1です（例題1）。もし光と同じように大きな角 $\alpha$ が使えたら、原子1個よりずっと小さい物まで見分けられるはずです。

ところが実際には、電子を曲げるレンズ（コイルの磁場で電子の向きを変える）は、ガラスのレンズほどきれいに像を結べず、レンズのふちを通った電子ほど、像からずれてしまいます（\textbf{収差}）。\bkidx{しゅうさ}{収差}そのため、角 $\alpha$ は 0.01 rad（約 0.6°）ほどしか使えません。それでも、200 kV の電子では $\delta\approx0.61\times\qty{2.5}{pm}\div0.01\approx\qty{0.15}{nm}$ で（例題12）、原子1個（約 0.1〜0.3 nm）を見分けられる細かさです。

\begin{bkfigure}
\begin{tikzpicture}[x=1mm, y=1mm]
  \draw[bkInk, line width=0.7pt, -{Stealth[length=2mm]}] (0,0) -- (128,0);
  \foreach \e/\t in {0/{1 mm}, 17/{\qty{100}{\micro m}}, 34/{\qty{10}{\micro m}}, 51/{\qty{1}{\micro m}}, 68/{\qty{100}{nm}}, 85/{\qty{10}{nm}}, 102/{\qty{1}{nm}}, 119/{0.1 nm}}{
    \draw[bkInk, line width=0.6pt] (\e+2,-1.2) -- (\e+2,1.2);
    \node[txd lab, anchor=north] at (\e+2,-1.8) {\t};}
  \foreach \x/\y/\t in {2/8/砂つぶ, 19/15/髪の毛, 36/8/細胞, 53/15/細菌, 70/8/ウイルス, 98.9/15/DNA のはば, 115.9/8/原子}{
    \draw[bkGray, line width=0.4pt] (\x,1.5) -- (\x,\y-0.5);
    \node[txd lab, anchor=south] at (\x,\y) {\t};}
  \fill[bkBlue!15] (2,-12) rectangle (64.9,-8);
  \node[txd lab, anchor=west] at (3,-10) {光の顕微鏡（約\qty{0.2}{\micro m}まで）};
  \fill[bkOrange!20] (2,-18) rectangle (121,-14);
  \node[txd lab, anchor=west] at (3,-16) {電子顕微鏡（約 0.1 nm：原子の大きさまで）};
\end{tikzpicture}
\bkfigcaption{物の大きさ（目もりは10倍ずつ）と、見分けられる細かさ。帯の右はしが、見分けられる限界。}\label{ph:原-3:fig-scale}
\end{bkfigure}

ドイツのルスカは、1931年に、コイルのレンズで電子の像を結ぶ、最初の電子顕微鏡を作りました。ルスカは翌年、「75 kV の電子と、開き角 0.02 rad を使えば、0.22 nm まで見分けられるはずだ」と見こみを立てました。その値が実際に達成されたのは、約40年後のことでした（ルスカのノーベル賞講演）。ルスカは1986年にノーベル物理学賞を受けました。

\begin{phmijika}[電子顕微鏡で見ているもの]
\textbf{透過電子顕微鏡（TEM）}は、うすくした試料に電子を通し、通りぬけた電子で像を作ります。原子の列や、半導体のチップの断面のトランジスタ（原-10）を写せます。\textbf{走査電子顕微鏡（SEM）}は、細くしぼった電子を試料の表面になぞるように当て、はね返ったり飛び出したりした電子を数えて、表面の立体的な像を作ります。ウイルスや虫の目、髪の毛の表面の写真の多くは SEM で撮ったものです。
\end{phmijika}

\Needspace{12\baselineskip}\bksubsection{7}{波束と不確定性：位置と運動量は同時に決まらない}
\label{ph:原-3:sec7}
②の平面波は、波長 $\lambda$ の波がどこまでも続いています。波長が1つに決まっているので、運動量 $p=h/\lambda$ はぴったり決まっています。けれども、波がどこにでも同じ高さで広がっているので、電子がどこにいるかは、まったくわかりません。

電子を「だいたいこのあたり」に置くには、波長が少しずつ違う波を重ねて、1か所だけ大きく、ほかでは打ち消し合う形を作ります。波長が少し違う2つの音を重ねると「うなり」ができ（4-9）、たくさん重ねると、1つのかたまりになります（9-3）。このかたまりを\textbf{波束}といいます。\bkidx{はそく}{波束}

\begin{bkfigure}
\begin{tikzpicture}
\begin{axis}[bkplot, name=a1, width=60mm, height=32mm, xmin=-30, xmax=30, ymin=-1.15, ymax=1.15, xtick=\empty, ytick={0}, xlabel={$x$}, title={\footnotesize 2つの波（うなり）}]
  \addplot[bkcurve, domain=-30:30, samples=400] {0.5*(cos(deg(1.0*x))+cos(deg(1.2*x)))};
\end{axis}
\begin{axis}[bkplot, at={(a1.east)}, anchor=west, xshift=12mm, width=60mm, height=32mm, xmin=-30, xmax=30, ymin=-1.15, ymax=1.15, xtick=\empty, ytick={0}, xlabel={$x$}, title={\footnotesize 7つの波}]
  \addplot[bkcurve, domain=-30:30, samples=400] {(cos(deg(0.85*x))+cos(deg(0.9*x))+cos(deg(0.95*x))+cos(deg(1.0*x))+cos(deg(1.05*x))+cos(deg(1.1*x))+cos(deg(1.15*x)))/7};
\end{axis}
\end{tikzpicture}
\begin{tikzpicture}
\begin{axis}[bkplot, width=124mm, height=34mm, xmin=-30, xmax=30, ymin=-1.15, ymax=1.15, xtick=\empty, ytick={0}, xlabel={$x$}, title={\footnotesize たくさんの波を、ガウスの形の重みで重ねた波束}]
  \addplot[bkcurve, domain=-30:30, samples=500] {exp(-x^2/36)*cos(deg(1.0*x))};
  \addplot[bkthin, dashed, domain=-30:30, samples=100] {exp(-x^2/36)};
  \addplot[bkthin, dashed, domain=-30:30, samples=100] {-exp(-x^2/36)};
\end{axis}
\end{tikzpicture}
\bkfigcaption{波長の少し違う波を重ねる。「2つの波」ではうなりが繰り返すが、「7つの波」では、かたまりとかたまりの間が広がる。たくさん重ねると、1つだけのかたまり（波束）になる（破線は、波束を包む形）。}\label{ph:原-3:fig-hasoku}
\end{bkfigure}

\Needspace{10\baselineskip}\noindent\textbf{表計算で波束を作る}\par\nobreak
「7つの波」のグラフは、表計算で作れます（9-3 の「いくつもの $\cos$ をたす」と同じ）。
\begin{enumerate}[label={\bkcircled[bkGreen]{\arabic*}}, leftmargin=2.4\zw, itemsep=1pt, topsep=2pt]
\item B〜H列の1行めに、7つの波数 0.85、0.90、…、1.15 を入れる。
\item A列の2行めから下に、位置 $x$ を 0、5、10、… と入れる（グラフをなめらかにかくときは、$-30$ から 30 まで 0.5 ずつ）。
\item B2 に \texttt{=COS(B\$1*\$A2)} と入れて、H列まで右へ、下へコピーする。I2 に7つの平均 \texttt{=AVERAGE(B2:H2)} を入れて下へコピーし、A列とI列でグラフをかく。
\end{enumerate}
\begin{bkcenter}
\bksheet{9}{10.8}{%
1 & $x$ ＼ $k$ & 0.85 & 0.90 & 0.95 & 1.00 & 1.05 & 1.10 & 1.15 & 平均 \\
2 & 0 & 1.000 & 1.000 & 1.000 & 1.000 & 1.000 & 1.000 & 1.000 & 1.000 \\
3 & 5 & $-0.446$ & $-0.211$ & 0.038 & 0.284 & 0.512 & 0.709 & 0.861 & 0.249 \\
4 & 10 & $-0.602$ & $-0.911$ & $-0.997$ & $-0.839$ & $-0.476$ & 0.004 & 0.483 & $-0.477$ \\
5 & 15 & 0.983 & 0.595 & $-0.113$ & $-0.760$ & $-0.999$ & $-0.702$ & $-0.029$ & $-0.146$ \\
6 & 20 & $-0.275$ & 0.660 & 0.989 & 0.408 & $-0.548$ & $-1.000$ & $-0.533$ & $-0.043$ \\}
\bktabcaption{波束の表計算。$x=0$ では7つの波の山がそろって、平均は1。はなれると、山と谷が打ち消し合って0に近づく。}\label{ph:原-3:tab-sheet}
\end{bkcenter}
重ねる波の数を増やし、波数の間かくをせまくするほど、次にまた山がそろう所（7つの波では $x=2\pi/0.05\approx126$）は遠くなり、1つだけのかたまりに近づきます。波数の広がり（0.85〜1.15）を半分にすると、かたまりのはばはほぼ2倍になります。これが次の話の出発点です。

\noindent\textbf{かたまりがせまいほど、たくさんの波長がいる}\par\nobreak
数学の本（9-4⑨）で、ガウスの形のパルスを波に分けると、時間のはばが $\sigma$ なら周波数のはばは $1/\sigma$ になり、「短いパルスほど、たくさんの周波数をふくむ」ことを確かめました。時間 $t$ を位置 $x$ に、角周波数 $\omega$ を波数 $k=2\pi/\lambda$ にかえれば、そのまま波束の話になります。
\begin{itemize}[leftmargin=1.2\zw, itemsep=2pt, topsep=2pt, label={\color{bkTheme}●}]
\item 位置のはばがせまい波束 → 波数（波長）のはばが広い。
\item 波数のはばがせまい波 → 位置のはばが広い（平面波は、波数が1つで、位置はどこまでも広がる）。
\end{itemize}

\Needspace{6\baselineskip}\noindent\textbf{はばのかけ算は $\tfrac12$ より小さくならない}\par\nobreak
電子の見つかる確率は $|\psi|^2$ でした（⑤）。そこで、$|\psi|^2$ のばらつきを、位置の\textbf{不確かさ} $\Delta x$（標準偏差。3-17）とします。波の高さが $\psi\propto e^{-x^2/(4\Delta x^2)}$ のガウスの形なら、$|\psi|^2\propto e^{-x^2/(2\Delta x^2)}$ で、はばがちょうど $\Delta x$ の正規分布です。
\begin{enumerate}[label={\bkcircled[bkGreen]{\arabic*}}, leftmargin=2.4\zw, itemsep=1pt, topsep=2pt]
\item 波の高さ $\psi$ は、$e^{-x^2/(2\sigma^2)}$ で $\sigma=\sqrt2\,\Delta x$ の形。9-4⑨から、波数の世界の高さのはばは $1/\sigma=1/(\sqrt2\,\Delta x)$。
\item 確率は高さの2乗。ガウスの形を2乗すると、はばは $1/\sqrt2$ 倍になる（$e^{-y^2/(2s^2)}$ の2乗は $e^{-y^2/(2(s/\sqrt2)^2)}$）。
\item 位置の確率のはばは $\sigma/\sqrt2=\Delta x$、波数の確率のはばは $\Delta k=\dfrac{1}{\sqrt2\,\Delta x}\times\dfrac{1}{\sqrt2}=\dfrac{1}{2\Delta x}$。
\end{enumerate}
だから、ガウスの形の波束では $\Delta x\,\Delta k=\tfrac12$ です（例題10）。ガウスでないほかの形では、かけ算はもっと大きくなることが、数学で確かめられています。つまり、どんな波でも
\[ \Delta x\,\Delta k\geqq\frac12 \]
です。最後に、$p=h/\lambda=\dfrac{h}{2\pi}\times\dfrac{2\pi}{\lambda}=\hbar k$ を使います。$\hbar=h/(2\pi)=\qty{1.055e-34}{J.s}$（エイチバー）は、$2\pi$ がよく出てくるので作った定数で、\textbf{ディラック定数}ともいいます。\bkidx{でぃらっくていすう}{ディラック定数}$\Delta p=\hbar\,\Delta k$ なので、
\pheq{\Delta x\,\Delta p\geqq\frac{\hbar}{2}}{m × kg·m/s ＝ J·s}
\label{ph:原-3:fukakutei}これを\textbf{不確定性関係}といいます。\bkidx{ふかくていせいかんけい}{不確定性関係}1927年に、ドイツのハイゼンベルクが示しました（1932年ノーベル物理学賞）。

\begin{phyomi}[$\Delta x\,\Delta p\geqq\dfrac{\hbar}{2}$]
「位置の不確かさと運動量の不確かさのかけ算は、$\hbar$ の半分より小さくできない」と読みます。位置をせまく決めようとすると、運動量（速さ）が広くばらつき、運動量をそろえようとすると、位置が広がります。これは、はかる道具が悪いからではありません。\textbf{電子が波であるために、もともと、両方が同時にぴったり決まった状態がない}のです。
\end{phyomi}

\begin{phbaisuru}[半分]{閉じこめる場所のはば $\Delta x$}
\noindent\begin{minipage}[c]{0.44\linewidth}\raggedright
運動量のばらつき $\Delta p$ は2倍以上になります。運動量のばらつきがあると、電子は止まっていられず、運動エネルギーの目安 $(\Delta p)^2/(2m)$ は4倍になります。0.1 nm の場所に閉じこめた電子は約 1 eV、0.05 nm なら約 4 eV です（例題7）。\textbf{せまい所に閉じこめるほど、電子は大きなエネルギーをもつ}のです。
\end{minipage}\hfill
\begin{minipage}[c]{0.54\linewidth}\centering
\begin{tikzpicture}
\begin{axis}[bkplot, name=b1, width=34mm, height=30mm, xmin=-3, xmax=3, ymin=0, ymax=1.15, xtick={0}, ytick=\empty, xlabel={$x$}, title={\footnotesize 位置}]
  \addplot[bkcurve, domain=-3:3, samples=80] {exp(-x^2/2)};
  \addplot[bkcurve2, dashed, domain=-3:3, samples=80] {exp(-x^2/(2*0.25))};
\end{axis}
\begin{axis}[bkplot, at={(b1.east)}, anchor=west, xshift=8mm, width=34mm, height=30mm, xmin=-3, xmax=3, ymin=0, ymax=1.15, xtick={0}, ytick=\empty, xlabel={$p$}, title={\footnotesize 運動量}]
  \addplot[bkcurve, domain=-3:3, samples=80] {exp(-x^2/(2*0.25))};
  \addplot[bkcurve2, dashed, domain=-3:3, samples=80] {exp(-x^2/2)};
\end{axis}
\end{tikzpicture}
\par\footnotesize 実線：もとの波束　破線：はばを半分にした波束
\end{minipage}
\end{phbaisuru}

\begin{phmijika}[原子が、ある大きさをもつわけ]
原子の中の電子は、原子核に引かれて、近づくほど位置エネルギーが下がります。それなら、電子は原子核に落ちてしまいそうです。けれども、せまい所に閉じこめると、不確定性関係で運動エネルギーが大きくなります。位置エネルギーが下がる得と、運動エネルギーが上がる損がつり合う所が、原子の大きさ（約 0.1 nm）と、原子の中の電子のエネルギー（数 eV）を決めているのです（原-4 ⑤・原-5）。
\end{phmijika}

\begin{phflow}[波を重ねることから、不確定性関係へ]
\phflowrow{ガウスの形のはば\\$\sigma\leftrightarrow\dfrac1\sigma$（9-4）}{確率は高さの2乗\\はばは $1/\sqrt2$ 倍}{$\Delta x\,\Delta k\geqq\dfrac12$}
\phflowrow{$\Delta x\,\Delta k\geqq\dfrac12$}{$p=\hbar k$\\（$\lambda=h/p$）}{$\Delta x\,\Delta p\geqq\dfrac{\hbar}{2}$}
\end{phflow}

\begin{seikaku}
ハイゼンベルクの1927年の論文は、「電子の位置を光で見ると、光子にけられて運動量が変わる」という思考実験で、$\Delta x\,\Delta p$ がおよそ $h$ の大きさになることを示しました。いまの $\Delta x\,\Delta p\geqq\hbar/2$ の形（標準偏差で書いた形）は、そのすぐあとに、ほかの研究者によって、波の性質だけから証明されました。この章の導き方は、こちらの考え方です。
\end{seikaku}

\Needspace{10\baselineskip}\bksubsection{8}{エネルギーと時間：スペクトルの線のはば}
\label{ph:原-3:sec8}
同じことは、時間と周波数の間にもあります。数学の本（9-13）で、「短い時間だけ見ると、周波数ははっきりしない」というきまりを見ました。光の波のパルスが時間 $\Delta t$ しか続かなければ、その光の周波数は $\Delta f\approx1/(2\pi\Delta t)$ ほどのはばをもちます。光子のエネルギーは $E=hf$ なので、エネルギーのはばは $\Delta E=h\,\Delta f$ です。まとめると、
\pheq{\Delta E\,\Delta t\sim\hbar}{J × s ＝ J·s}
\label{ph:原-3:et}です（$\sim$ は「大きさがだいたい同じ」の意味）。

\noindent\begin{minipage}[c]{0.5\linewidth}\raggedright
原子の中で高いエネルギーの段にいる電子は、平均の寿命 $\tau$ ほどたつと光を出して下の段に移ります（原-4 ⑪）。光を出す時間がおよそ $\tau$ なので、出てくる光の周波数は、ぴったり1つの値ではなく、$\Delta f=1/(2\pi\tau)$ ほどのはばをもちます。これを、スペクトルの線の\textbf{自然幅}といいます。\bkidx{しぜんはば}{自然幅}寿命の短い段ほど、線は太くなります。
\end{minipage}\hfill
\begin{minipage}[c]{0.47\linewidth}\centering
\begin{tikzpicture}
\begin{axis}[bkplot, axis lines=left, width=56mm, height=38mm, xmin=-5, xmax=5, ymin=0, ymax=1.2, xtick=\empty, ytick={0.5,1}, xlabel={周波数 $f$}, ylabel={光の強さ}]
  \addplot[bkcurve, domain=-5:5, samples=200] {1/(1+4*x^2)};
  \addplot[bkcurve2, dashed, domain=-5:5, samples=200] {1/(1+(4*x^2)/9)};
  \draw[bkInk, line width=0.5pt, {Stealth[length=1.4mm]}-{Stealth[length=1.4mm]}] (axis cs:-0.5,0.5) -- (axis cs:0.5,0.5);
  \node[font=\sffamily\footnotesize, anchor=north] at (axis cs:0,0.45) {$\Delta f$};
\end{axis}
\end{tikzpicture}
\bkfigcaption{線の形。寿命を $\tfrac13$ にすると、はばは3倍に（破線）。}\label{ph:原-3:fig-haba}
\end{minipage}

ナトリウムの黄色い光（\qty{589}{nm}）を出す段の寿命は約 \qty{16}{ns} で、自然幅は約 \qty{10}{MHz} です（例題9）。光の周波数 \qty{5.1e14}{Hz} の約5000万分の1（$2\times10^{-8}$）しかありません。原子の光の色が、とても正確に決まっているのはこのためで、この性質が原子時計やレーザーの土台になっています。

\Needspace{12\baselineskip}\bksubsection{9}{群速度：電子の波の速さは、つぶの速さ}
\label{ph:原-3:sec9}
電子の波は、どんな速さで進むのでしょう。光子では $E=hf$ でした。電子の波にも同じ $E=hf=\hbar\omega$ を当てはめ（$\omega=2\pi f$）、$p=\hbar k$ と組み合わせます。光の速さより遅い電子のエネルギーは運動エネルギー $E=p^2/(2m)$ なので、
\pheq{\hbar\omega=\frac{(\hbar k)^2}{2m}\qquad\text{つまり}\qquad\omega=\frac{\hbar k^2}{2m}}{rad/s ＝ J·s × (rad/m)$^2$ ÷ kg}
\label{ph:原-3:bunsan}です。$\omega$ と $k$ の関係を\textbf{分散関係}といいました（波-5 ⑦）。音や弦の波では $\omega=ck$ で、$\omega$ は $k$ に比例しました。電子の波では、$\omega$ が $k$ の2乗に比例します。

山や谷が進む速さ（位相の速さ）は $\omega/k$ でした。
\[ \frac{\omega}{k}=\frac{\hbar k}{2m}=\frac{p}{2m}=\frac{v}{2} \]
で、電子の速さ $v=p/m$ の\textbf{半分}です。それでは、電子の波は、電子より遅れて進むのでしょうか。

そうではありません。電子は波束で、波束のかたまりは\textbf{群速度} $v_{\mathrm g}=\dd\omega/\dd k$ で進みます（波-5 ⑦。波-5 では $c_{\mathrm g}$ と書いた）。\bkidx{ぐんそくど}{群速度}
\pheq{v_{\mathrm g}=\frac{\dd\omega}{\dd k}=\frac{\hbar k}{m}=\frac{p}{m}=v}{m/s}
\label{ph:原-3:gun}波束のかたまりは、電子の速さ $v$ でちょうど進みます。電子が見つかる場所（$|\psi|^2$ の大きい所）は波束のかたまりなので、\textbf{電子は、つぶとしての速さで進む}のです。山と谷は、波束より遅いので、波束の中を前から後ろへ、ゆっくりずれていきます（図\ref{ph:原-3:fig-gun}）。

\begin{bkfigure}
\begin{tikzpicture}
\begin{axis}[bkplot, axis lines=none, width=124mm, height=46mm, xmin=-22, xmax=36, ymin=-1.75, ymax=1.65]
  \draw[bkGray, line width=0.5pt, -{Stealth[length=1.6mm]}] (axis cs:-22,0) -- (axis cs:36,0) node[anchor=west, font=\footnotesize] {$x$};
  \addplot[bkcurve, domain=-22:18, samples=500] {exp(-x^2/100)*cos(deg(1.5*x))};
  \addplot[bkcurve2, dashed, domain=-6:36, samples=500] {exp(-(x-12)^2/100)*cos(deg(1.5*x-9))};
  \draw[bkGray, line width=0.4pt, densely dotted] (axis cs:0,-1.3) -- (axis cs:0,1.45) (axis cs:12,-0.2) -- (axis cs:12,1.45) (axis cs:6,-1.3) -- (axis cs:6,-0.2);
  \draw[bkRed!80!black, line width=0.9pt, -{Stealth[length=2mm]}] (axis cs:0,1.45) -- (axis cs:12,1.45);
  \node[font=\sffamily\footnotesize, anchor=west] at (axis cs:12.5,1.45) {かたまりの中心は $vt$ 進む};
  \fill[bkBlue!80!black] (axis cs:-0.6,1.28) -- (axis cs:0.6,1.28) -- (axis cs:0,1.06) -- cycle;
  \fill[bkOrange!85!black] (axis cs:5.4,0.98) -- (axis cs:6.6,0.98) -- (axis cs:6,0.76) -- cycle;
  \draw[bkOrange!85!black, line width=0.9pt, -{Stealth[length=2mm]}] (axis cs:0,-1.3) -- (axis cs:6,-1.3);
  \node[font=\sffamily\footnotesize, anchor=west] at (axis cs:6.5,-1.3) {印の山は $\tfrac12vt$ しか進まない};
\end{axis}
\end{tikzpicture}
\bkfigcaption{電子の波束が進むようす。はじめの波束（実線）と、時間 $t$ のあと（破線）。かたまりの中心は $vt$ 進むが、はじめ中心にあった山（三角の印）は $vt/2$ しか進まない。山はかたまりの中を後ろへずれていき、新しい山が前から現れる。}\label{ph:原-3:fig-gun}
\end{bkfigure}

\begin{seikaku}
実際の電子の波束は、進むうちに、少しずつ広がっていきます。分散関係が $\omega\propto k$ でないので、波長の違う成分の進む速さが少しずつ違うからです（波-5 ⑦の分散）。また、位相の速さは、エネルギーの0をどこに選ぶかで変わります（相対論の全エネルギー $mc^2$ をふくめると、位相の速さは光の速さより大きくなる）。はかれるのは群速度のほうで、こちらはエネルギーの0の選び方によりません。
\end{seikaku}

\bksubsection{10}{中性子や原子も波}
\label{ph:原-3:sec10}
ド・ブロイ波長は、電子だけのものではありません。原子炉から出る中性子を、周りの温度（約20 ℃）のエネルギー 0.0253 eV まで遅くした\textbf{熱中性子}は、速さ約 \qty{2200}{m/s}、波長 \qty{0.18}{nm} で、原子の間かくと同じくらいです（例題5）。中性子は電気をもたないので、物の奥まで入り、結晶の中の原子の並びや、磁石の中の小さな磁石の向きを調べるのに使われます（原-13・原-15）。

もっと重い物も波としてふるまいます。1999年、オーストリアの研究グループは、炭素の原子60個がサッカーボールの形につながった分子（C$_{60}$）でも、干渉のしまを写しました。それより大きな分子でも、干渉が確かめられています。では、どこまで大きな物が波としてふるまうのか、ボールのような身の回りの物で波が見えないのはなぜか（例題4 では波長が短すぎる、というだけでよいのか）は、次のまだわかっていないことにつながる問題です。

\Needspace{16\baselineskip}\begin{motto}[4]
\phboxhead{サッカーボールの形の分子の干渉を計算する}
1999年の実験では、C$_{60}$ を約 900 K（約 630 ℃）の炉で気体にして飛ばし、溝の間かく \qty{100}{nm} の回折格子に通して、\qty{1}{m} 先で分子を1個ずつ数えました。報告された波長は \qty{2.8}{pm} です。C$_{60}$ の質量は $720\times\qty{1.66e-27}{kg}\approx\qty{1.20e-24}{kg}$ なので、速さは
\[ v=\frac{h}{m\lambda}=\frac{6.63\times10^{-34}}{1.20\times10^{-24}\times2.8\times10^{-12}}\approx\qty{200}{m/s} \]
です。1次の回折の向きは $\sin\theta=\lambda/d=2.8\times10^{-12}\div(100\times10^{-9})=\num{2.8e-5}$ で、1 m 先では、まっすぐの向きから約 \qty{28}{\micro m} しかはなれていません。分子の大きさ（約 1 nm）は、波長の約400倍もあります。それでも、分子1個の波が格子の何本ものすきまを同時に通って干渉したのです。
\end{motto}

%% ---------------------------------------------------------------
\begin{phmada}[「観測すると1点に決まる」とは、どういうことか]
\label{ph:原-3:sec11}
⑤の実験で、電子1個の波は、2つの道を同時に通ってきたように広がっています。それなのに、スクリーンにぶつかると、電子はいつも1つの点で見つかります。1個ずつの行き先は、$|\psi|^2$ で決まる確率でしか言えません。

「広がった波が、見つかった瞬間に1点に決まる」ことを、どう理解すればよいのでしょう。計算の規則（$|\psi|^2$ が確率を与える）は、これまでのどの実験とも合っています。けれども、その規則が「何を意味しているのか」――波は本当にあるものなのか、私たちの知識を表しているだけなのか、見つかる前に電子はどこにいたのか――については、約100年たった今も、物理学者の考えは1つにまとまっていません。この問題を\textbf{観測の問題}といいます。\bkidx{かんそくのもんだい}{観測の問題}

2つの電子や光子の状態が、はなれていても切りはなせない形でつながる「量子のもつれ」の実験は、身近な感覚では説明できない結びつきが本当にあることを確かめ、2022年のノーベル物理学賞になりました。この性質は、量子コンピュータや、こっそり読まれない通信（量子暗号）に使われ始めています。
\end{phmada}

\begin{phsaki}{量子力学と量子情報へ}
この章では、電子の波の形をガウスの形などと決めて考えました。波の形が、力（位置エネルギー）によってどう決まるのかを表す式が、シュレディンガー方程式です（原-5）。それを使うと、箱や原子に閉じこめた電子のエネルギーが、とびとびの値になることや、電子が壁をすりぬけること（原-6）が計算できます。不確定性関係は、専門の本では、位置と運動量を「演算子」という計算の記号で表し、その順番を入れかえると結果が変わることから導きます（原-5 の専門書をのぞいてみよう）。量子のもつれは、量子情報の科学として、専門の本で学びます。
\end{phsaki}

\begin{senmonsho}[量子力学・電子顕微鏡の専門書の書き方]
\begin{bkhonbun}
加速電圧 $E$ の電子の波長は，相対論補正を含めて
\[ \lambda=\frac{h}{\sqrt{2m_0eE\,(1+eE/2m_0c^2)}} \]
となる（補正前の値は $\lambda\simeq12.26/\sqrt{E}$ Å，$E$ は V 単位）．位置と運動量の標準偏差 $\sigma_x$，$\sigma_p$ の間には不確定性関係 $\sigma_x\sigma_p\geqq\hbar/2$ が成り立つ．
\end{bkhonbun}
\bkyaku{専門書では、加速電圧を $E$、止まっているときの電子の質量を $m_0$（静止質量）と書くことがある。Å（オングストローム）は 0.1 nm で、結晶や電子顕微鏡の本でよく使われる長さの単位。$\sigma_x$・$\sigma_p$ は、この章の $\Delta x$・$\Delta p$ と同じ（標準偏差）。$\simeq$ は「ほぼ等しい」。}
\end{senmonsho}

\begin{taisetsu}
\def\tsl#1{\makebox[7.4\zw][l]{\sffamily\bfseries\small\color{\bkTT}#1}}
\tsl{ド・ブロイ波長}$\lambda=\dfrac{h}{p}=\dfrac{h}{mv}$。電圧 $V$ で加速した電子は $\lambda=\dfrac{h}{\sqrt{2meV}}$\\[3pt]
\tsl{電子の回折}$a\sin\theta=n\lambda$（表面の原子の列）。結晶の面は $2d\sin\theta=n\lambda$（原-8）\\[3pt]
\tsl{確率の波}見つかる確率は $|\psi|^2$ に比例。1個ずつ届き、集まると干渉のしま\\[3pt]
\tsl{見分ける細かさ}$\delta\approx0.61\lambda/\alpha$。電子顕微鏡は約 0.1 nm\\[3pt]
\tsl{不確定性関係}$\Delta x\,\Delta p\geqq\dfrac{\hbar}{2}$、$\Delta E\,\Delta t\sim\hbar$。$\hbar=\dfrac{h}{2\pi}$、$p=\hbar k$\\[3pt]
\tsl{群速度}$\omega=\dfrac{\hbar k^2}{2m}$。$v_{\mathrm g}=\dfrac{\dd\omega}{\dd k}=v$（位相の速さは $v/2$）
\end{taisetsu}

\begin{kotoba}
\small
\bkword{ド・ブロイ波長}{つぶの運動量 $p$ から決まる波長 $h/p$。}
\bkword{物質波}{電子・中性子・原子などの、つぶがもつ波。ド・ブロイ波ともいう。}
\bkword{電子線回折}{結晶に電子を当てたとき、決まった向きに電子が強め合って進むこと。}
\bkword{波束}{波長の少しずつ違う波を重ねて、1か所に集めた波のかたまり。}
\bkword{不確定性関係}{位置と運動量の不確かさのかけ算が $\hbar/2$ より小さくならないこと。}
\bkword{群速度}{波束のかたまりが進む速さ $\dd\omega/\dd k$。電子では電子の速さと同じ。}
\end{kotoba}

\bkkigocards{\bkkigocard{$\hbar$}{エイチ・バー}{ディラック定数 $h/(2\pi)$〔J·s〕。$p=\hbar k$、$E=\hbar\omega$}{reduced Planck constant}{原-3}}
{\bkkigocard{$\Delta x$}{デルタ・エックス}{位置の不確かさ（ばらつきの標準偏差）。運動量の不確かさは $\Delta p$}{uncertainty}{原-3}}
\bkkigocards{\bkkigocard{$\psi$}{プサイ}{電子などの波の高さ（波動関数）。見つかる確率は $|\psi|^2$ に比例（原-5 で式を立てる。時間をふくむと $\Psi$）}{wave function}{原-3}}
{\bkkigocard{$v_{\mathrm g}$}{ブイ・ジー}{群速度 $\dd\omega/\dd k$〔m/s〕。波-5 では $c_{\mathrm g}$ と書いた}{group velocity}{原-3}}

%% ======================== 例題 ========================
\Needspace{10\baselineskip}\bksection{reidai}{例題でたしかめよう}
\label{ph:原-3:ex1}
\begin{phreidai}[電子顕微鏡の電子の波長（はてな？）]
透過電子顕微鏡で、電子を 100 kV（10万 V）で加速します。電子のド・ブロイ波長を求め、緑の光（\qty{550}{nm}）の何分の1かを答えましょう。まず相対論を考えずに求め、次に相対論の直し（②のせいかくに言うと）を入れて、どれだけ変わるかを確かめます（3けたで書きます）。
\phbamen
電圧で加速された電子を、波として見る場面です（②⑥）。
\phmushi
加速する前の電子のエネルギーは0とします。電子顕微鏡の中の空気（真空にしてある）や、ほかの電子との力は無視します。
\phzu
$V=\qty{1.00e5}{V}$、$m=\qty{9.11e-31}{kg}$、$e=\qty{1.60e-19}{C}$、$h=\qty{6.63e-34}{J.s}$、$mc^2=\qty{511}{keV}$。
\phriyu
加速電圧から運動量がわかり、$\lambda=h/p$（②）で波長が出ます。100 kV は $mc^2$ の約5分の1なので、相対論の直しを確かめます。
\phtateru
\[ \lambda_0=\frac{h}{\sqrt{2meV}},\qquad \lambda=\frac{\lambda_0}{\sqrt{1+eV/(2mc^2)}} \]
\phtoku
\[ \lambda_0=\frac{1.226}{\sqrt{1.00\times10^5}}\ \mathrm{nm}=\qty{3.88e-3}{nm}=\qty{3.88}{pm} \]
\[ 1+\frac{eV}{2mc^2}=1+\frac{100}{2\times511}=1.098,\qquad \lambda=\frac{3.88}{\sqrt{1.098}}=\qty{3.70}{pm} \]
\[ \frac{\qty{550}{nm}}{\qty{3.70}{pm}}=\frac{550\times10^{-9}}{3.70\times10^{-12}}\approx\num{1.49e5} \]
\phtashikame
単位：J·s÷$\sqrt{\text{kg·J}}$＝$\sqrt{\text{J·s}^2/\text{kg}}$＝$\sqrt{\text{m}^2}$＝m。\\
大きさ：表\ref{ph:原-3:tab-lambda}と同じで、緑の光の約15万分の1。相対論の直しは約5\,\%です。\\
極端な場合：電圧を下げると、直しは小さくなります（1 kV で 0.05\,\%）。
\kotae 相対論なしで \qty{3.88}{pm}、直すと \qty{3.70}{pm}。緑の光の約15万分の1
\end{phreidai}
\begin{jibun}
電子を 50 kV で加速するときの波長を、相対論なしと、直したときで求めましょう（3けたで書く）。緑の光（\qty{550}{nm}）の約何万分の1ですか。
\bkjans{相対論なし \qty{5.48}{pm}。$1+50/1022=1.049$、その平方根でわって \qty{5.36}{pm}。約10万分の1}
\end{jibun}

\label{ph:原-3:ex2}
\begin{phreidai}[100 V の電子の波長]
100 V で加速した電子の速さと、ド・ブロイ波長を求めましょう。その波長は、原子の間かく（約 0.1〜0.3 nm）とくらべて、どれくらいですか。
\phbamen
ブラウン管や、デビソンとガーマーの実験と同じくらいの、遅い電子の波です（②）。
\phmushi
相対論は考えません（100 V は $mc^2$ の5千分の1）。
\phzu
$V=\qty{100}{V}$。運動エネルギー $eV$、速さ $v$、運動量 $p=mv$。
\phriyu
$\tfrac12mv^2=eV$ で速さが出て、$p=mv$ を $\lambda=h/p$ に入れます。
\phtateru
\[ v=\sqrt{\frac{2eV}{m}},\qquad \lambda=\frac{h}{mv} \]
\phtoku
\[ v=\sqrt{\frac{2\times1.60\times10^{-19}\times100}{9.11\times10^{-31}}}=\qty{5.93e6}{m/s} \]
\[ \lambda=\frac{6.63\times10^{-34}}{9.11\times10^{-31}\times5.93\times10^{6}}=\qty{1.23e-10}{m}=\qty{0.123}{nm} \]
\phtashikame
単位：J·s÷(kg·m/s)＝(kg·m$^2$/s)÷(kg·m/s)＝m。\\
大きさ：$1.226/\sqrt{100}=0.1226$ nm と一致。原子の間かくと同じ大きさなので、結晶で回折します（③）。\\
極端な場合：速さは光の速さの約2\,\% で、相対論を考えなくてよいことも確かめられます。
\kotae 速さ約\qty{5.9e6}{m/s}、波長約\qty{0.12}{nm}（原子の間かくと同じくらい）
\end{phreidai}
\begin{jibun}
ド・ブロイ波長が \qty{0.10}{nm} の電子の運動エネルギーは、何 eV ですか。何 V で加速すればよいですか。
\bkjans{$p=h/\lambda=\qty{6.63e-24}{kg.m/s}$。$p^2/(2m)\approx\qty{2.4e-17}{J}$、約\qty{150}{eV}。約 150 V}
\end{jibun}

\label{ph:原-3:ex3}
\begin{phreidai}[デビソンとガーマーの実験]
ニッケルの結晶の表面に、54 V で加速した電子を垂直に当てます。表面の原子の列の間かくを \qty{0.215}{nm} として、はね返る電子が強め合う向き（$n=1$）を求め、実験の 50° とくらべましょう（3けたで書きます）。
\phbamen
表面の原子の列を、反射型の回折格子とみなす場面です（③）。
\phmushi
結晶の奥の層ではね返る電子と、電子が結晶に入るときに速くなる効果は無視し、表面の1層だけで考えます。
\phzu
図\ref{ph:原-3:fig-dg}。$a=\qty{0.215}{nm}$、$V=\qty{54}{V}$、角 $\theta$ は表面に垂直な向きから。
\phriyu
となりの列との道のりの差 $a\sin\theta$ が1波長のとき強め合います（③）。波長は②の式。
\phtateru
\[ \lambda=\frac{1.226}{\sqrt{54}}\ \mathrm{nm},\qquad \sin\theta=\frac{\lambda}{a} \]
\phtoku
\[ \lambda=\frac{1.226}{7.348}=\qty{0.167}{nm},\qquad \sin\theta=\frac{0.167}{0.215}=0.776,\qquad\theta\approx\ang{50.9} \]
\phtashikame
単位：nm÷nm で、$\sin$ は単位なし。\\
大きさ：実験の 50° と約 1° の違いで合います。50° から逆に求めた波長 \qty{0.165}{nm} と、計算の \qty{0.167}{nm} は約1\,\% の違いです。\\
極端な場合：電圧を下げて $\lambda>a$ にすると、$\sin\theta>1$ になり、強め合う向きはありません（約 33 V より下）。
\kotae 約 51°（実験の 50° とよく合う）
\end{phreidai}
\begin{jibun}
例題3 で、電圧を 65 V にすると、強め合う向きは何度になりますか。
\bkjans{$\lambda=1.226/\sqrt{65}\approx\qty{0.152}{nm}$、$\sin\theta=0.152/0.215\approx0.708$ で、約 45°}
\end{jibun}

\label{ph:原-3:ex4}
\begin{phreidai}[野球のボールの波長]
質量 \qty{0.145}{kg} の野球のボールが \qty{40}{m/s}（時速 144 km）で飛んでいます。ド・ブロイ波長を求めましょう。ボールが回折しないのはなぜですか。
\phbamen
身の回りの物にも $\lambda=h/p$ を当てはめてみる場面です（②）。
\phmushi
ボールの回転や空気の抵抗は無視します。ボールを1つのつぶとします。
\phzu
$m=\qty{0.145}{kg}$、$v=\qty{40}{m/s}$。
\phriyu
光の速さよりずっと遅いので、$p=mv$ として $\lambda=h/(mv)$ が使えます。
\phtateru
\[ \lambda=\frac{h}{mv} \]
\phtoku
\[ \lambda=\frac{6.63\times10^{-34}}{0.145\times40}=\frac{6.63\times10^{-34}}{5.8}\approx\qty{1.1e-34}{m} \]
\phtashikame
単位：m。\\
大きさ：原子核（約 $10^{-15}$ m）よりも、さらに $10^{19}$ 倍も短い波長です。回折が目立つのは、波長と同じくらいのすきまや物の間かくがあるときでした（やってみよう）。そんなに短い間かくの物はないので、ボールの波としてのふるまいは、どんな実験でも見えません。\\
極端な場合：$m$ が電子ほど小さければ、波長は原子の大きさになります（例題2）。
\kotae 約\qty{1.1e-34}{m}。短すぎて、回折を起こすような間かくの物がない
\end{phreidai}
\begin{jibun}
質量 \qty{70}{kg} の人が \qty{5.0}{m/s} で走っているときの、ド・ブロイ波長を求めましょう。
\bkjans{$p=70\times5.0=\qty{350}{kg.m/s}$。$\lambda=h/p\approx\qty{1.9e-36}{m}$}
\end{jibun}

\label{ph:原-3:ex5}
\begin{phreidai}[熱中性子の速さと波長]
運動エネルギーが \qty{0.0253}{eV} の中性子（熱中性子）の速さと波長を求めましょう。中性子の質量は \qty{1.675e-27}{kg} です（3けたで書きます）。
\phbamen
原子炉の中で、周りの温度まで遅くなった中性子を、波として見る場面です（⑩）。
\phmushi
中性子の向きのばらつき（いろいろなエネルギーの中性子がまざっていること）は無視し、1つのエネルギーで考えます。
\phzu
$K=\qty{0.0253}{eV}=0.0253\times\qty{1.602e-19}{J}$、$m_{\mathrm n}=\qty{1.675e-27}{kg}$。
\phriyu
$K=\tfrac12m_{\mathrm n}v^2$ から速さ、$\lambda=h/(m_{\mathrm n}v)$ から波長が出ます。
\phtateru
\[ v=\sqrt{\frac{2K}{m_{\mathrm n}}},\qquad\lambda=\frac{h}{m_{\mathrm n}v} \]
\phtoku
\[ v=\sqrt{\frac{2\times0.0253\times1.602\times10^{-19}}{1.675\times10^{-27}}}=\sqrt{\num{4.84e6}}\approx\qty{2.20e3}{m/s} \]
\[ \lambda=\frac{6.63\times10^{-34}}{1.675\times10^{-27}\times2200}\approx\qty{1.80e-10}{m}=\qty{0.180}{nm} \]
\phtashikame
単位：m/s と m。\\
大きさ：熱中性子の速さは「2200 m/s」と決まった値でよばれ、それと合います。波長は原子の間かくと同じくらいです。\\
極端な場合：同じエネルギーの電子なら、質量が1839分の1なので、波長は $\sqrt{1839}\approx43$ 倍長くなります。
\kotae 約\qty{2.20e3}{m/s}、約\qty{0.180}{nm}
\end{phreidai}
\begin{jibun}
100 V で加速した陽子（質量 \qty{1.67e-27}{kg}）のド・ブロイ波長を求めましょう。同じ電圧の電子の何分の1ですか。
\bkjans{$h/\sqrt{2m_{\mathrm p}eV}$ で約 \qty{2.9}{pm}。電子の $\sqrt{m_{\mathrm p}/m}\approx43$ 分の1}
\end{jibun}

\label{ph:原-3:ex6}
\begin{phreidai}[電子の二重スリットのしま]
50 kV で加速した電子を、間かく \qty{1.0}{\micro m} の2つのすきまに通し、\qty{1.0}{m} はなれたスクリーンに当てます。しまの間かくを求めましょう。波長は相対論で直した \qty{5.36}{pm}（例題1 のじぶんでとこう）を使います。
\phbamen
光の二重スリット（波-13 ②）の光を、電子の波にかえた場面です（⑤）。
\phmushi
すきまの幅の効果（全体がゆるやかに暗くなる）は無視し、2つの道の波の重ね合わせだけを考えます。電子は1個ずつでもよい。
\phzu
$d=\qty{1.0e-6}{m}$、$L=\qty{1.0}{m}$、$\lambda=\qty{5.36e-12}{m}$。
\phriyu
すきまの間かくより、スクリーンまでの距離がずっと大きいので、波-13 ②の $\Delta x=L\lambda/d$ が使えます。
\phtateru
\[ \Delta x=\frac{L\lambda}{d} \]
\phtoku
\[ \Delta x=\frac{1.0\times5.36\times10^{-12}}{1.0\times10^{-6}}\approx\qty{5.4e-6}{m}=\qty{5.4}{\micro m} \]
\phtashikame
単位：m×m÷m＝m。\\
大きさ：髪の毛の太さの約20分の1と、とても細かいしまです。実際の実験では、電子のレンズで像を何千倍にも大きくして記録します。\\
極端な場合：電圧を上げて波長を短くすると、しまはさらに細かくなります。
\kotae 約\qty{5.4}{\micro m}
\end{phreidai}
\begin{jibun}
100 kV の電子（波長 \qty{3.70}{pm}）を、間かく \qty{0.50}{\micro m} の2つのすきまに通し、\qty{2.0}{m} 先のスクリーンに当てます。しまの間かくを求めましょう。
\bkjans{$\Delta x=L\lambda/d$ で、約\qty{15}{\micro m}}
\end{jibun}

\label{ph:原-3:ex7}
\begin{phreidai}[原子に閉じこめた電子のエネルギー]
電子を、原子の大きさくらいの $\Delta x=\qty{0.10}{nm}$ の中に閉じこめると、運動量と速さのばらつきは少なくともいくらになりますか。そのばらつきから、運動エネルギーの大きさの目安 $(\Delta p)^2/(2m)$ を eV で求めましょう。
\phbamen
原子の中の電子のエネルギーの大きさを、不確定性関係で見積もる場面です（⑦）。
\phmushi
波束の細かい形（ガウスかどうか）は気にしません。大きさの目安（けた）がわかればよいとします。
\phzu
$\Delta x=\qty{1.0e-10}{m}$、$\hbar=\qty{1.055e-34}{J.s}$、$m=\qty{9.11e-31}{kg}$。
\phriyu
不確定性関係 $\Delta x\,\Delta p\geqq\hbar/2$ の等号で、いちばん小さいばらつきが出ます。運動量がばらつくと、平均の運動量が0でも、運動エネルギーは $(\Delta p)^2/(2m)$ ほどあります。
\phtateru
\[ \Delta p=\frac{\hbar}{2\Delta x},\qquad\Delta v=\frac{\Delta p}{m},\qquad K\approx\frac{(\Delta p)^2}{2m} \]
\phtoku
\[ \Delta p=\frac{1.055\times10^{-34}}{2\times1.0\times10^{-10}}\approx\qty{5.3e-25}{kg.m/s} \]
\[ \Delta v=\frac{5.3\times10^{-25}}{9.11\times10^{-31}}\approx\qty{5.8e5}{m/s} \]
\[ K\approx\frac{(5.3\times10^{-25})^2}{2\times9.11\times10^{-31}}\approx\qty{1.5e-19}{J}\approx\qty{1}{eV} \]
\phtashikame
単位：J·s÷m＝kg·m/s。\\
大きさ：原子の中の電子のエネルギーは数 eV（水素では 13.6 eV：原-4 ⑤）で、目安として合っています。速さのばらつきは秒速約 600 km です。\\
極端な場合：$\Delta x$ を半分にすると、$K$ は4倍の約 4 eV になります。
\kotae $\Delta p\approx\qty{5.3e-25}{kg.m/s}$、$\Delta v\approx\qty{5.8e5}{m/s}$、約\qty{1}{eV}
\end{phreidai}
\begin{jibun}
電子を、はば $\Delta x=\qty{1.0}{nm}$ の小さな半導体のつぶ（量子ドット）に閉じこめたときの、$(\Delta p)^2/(2m)$ の目安を eV で求めましょう。
\bkjans{0.1 nm のときの $(1/10)^2$ 倍で、約 \qty{0.01}{eV}（\qty{9.5e-3}{eV}）}
\end{jibun}

\label{ph:原-3:ex8}
\begin{phreidai}[ちりの不確定性は気にしなくてよい]
質量 \qty{1}{\micro g}（\qty{1e-9}{kg}）のちりを、\qty{1}{\micro m} の場所に閉じこめたとします。速さのばらつきは少なくともいくらですか。
\phbamen
身の回りの小さな物に、不確定性関係を当てはめる場面です（⑦）。
\phmushi
空気の分子にぶつかってゆれる動き（ブラウン運動）は考えません。
\phzu
$m=\qty{1e-9}{kg}$、$\Delta x=\qty{1e-6}{m}$。
\phriyu
$\Delta x\,\Delta p\geqq\hbar/2$ から、$\Delta v=\Delta p/m$ のいちばん小さい値を出します。
\phtateru
\[ \Delta v=\frac{\hbar}{2m\,\Delta x} \]
\phtoku
\[ \Delta v=\frac{1.055\times10^{-34}}{2\times10^{-9}\times10^{-6}}\approx\qty{5e-20}{m/s} \]
\phtashikame
単位：J·s÷(kg·m)＝m/s。\\
大きさ：1年（約 \qty{3e7}{s}）たっても、動く距離は $10^{-12}$ m ほどで、原子の大きさの100分の1です。ちりの位置と速さは、ふつうの意味で両方決まっていると考えてよいのです。\\
極端な場合：質量を電子にすると、速さのばらつきは $10^{21}$ 倍になります。
\kotae 約\qty{5e-20}{m/s}（気にしなくてよい）
\end{phreidai}
\begin{jibun}
電子を \qty{1}{\micro m} の場所に閉じこめたときの、速さのばらつきの最小の値を求めましょう。
\bkjans{$\Delta v=\hbar/(2m\,\Delta x)$ で、約\qty{58}{m/s}}
\end{jibun}

\label{ph:原-3:ex9}
\begin{phreidai}[ナトリウムの光の線のはば]
ナトリウムの黄色い光（\qty{589}{nm}）を出す段は、1秒あたり \num{6.16e7} の割合で光を出して下の段に移ります（平均の寿命は、その逆数）。光の線の自然幅 $\Delta f=1/(2\pi\tau)$ を求め、光の周波数とくらべましょう。
\phbamen
光を出す時間が限られているために、光の周波数がぼやける場面です（⑧）。
\phmushi
原子が動くことによるドップラー効果のはば（波-8）や、原子どうしがぶつかることによるはばは無視し、寿命によるはばだけを考えます。
\phzu
$\tau=1/(\num{6.16e7}\,\text{/s})$、$\lambda=\qty{589e-9}{m}$。
\phriyu
光を出す時間が $\tau$ くらいなので、周波数のはばは $1/(2\pi\tau)$ ほど（⑧）。光の周波数は $f=c/\lambda$。
\phtateru
\[ \tau=\frac{1}{\num{6.16e7}},\qquad\Delta f=\frac{1}{2\pi\tau},\qquad f=\frac{c}{\lambda} \]
\phtoku
\[ \tau\approx\qty{1.62e-8}{s}=\qty{16.2}{ns},\qquad\Delta f=\frac{1}{2\pi\times1.62\times10^{-8}}\approx\qty{9.8e6}{Hz} \]
\[ f=\frac{3.00\times10^8}{589\times10^{-9}}\approx\qty{5.09e14}{Hz},\qquad\frac{\Delta f}{f}\approx\num{1.9e-8} \]
\phtashikame
単位：1/s＝Hz。\\
大きさ：約 \qty{10}{MHz} のはばは、光の周波数の約5000万分の1です。線はとても細いのです。\\
極端な場合：寿命が長い段ほど、線は細くなります。原子時計では、寿命の長い段の、とても細い線を使います。
\kotae $\tau\approx\qty{16}{ns}$、$\Delta f\approx\qty{9.8}{MHz}$（光の周波数の約 $2\times10^{-8}$）
\end{phreidai}
\begin{jibun}
寿命が \qty{1.0}{\micro s} の段から出る光の自然幅は、何 MHz ですか。
\bkjans{$\Delta f=1/(2\pi\tau)$ で約\qty{0.16}{MHz}}
\end{jibun}

\label{ph:原-3:ex10}
\begin{phreidai}[ガウスの形の波束]
電子の見つかる確率 $|\psi|^2$ が、はば（標準偏差）$\Delta x=\qty{1.0}{nm}$ のガウスの形の波束があります。波数のはば $\Delta k$ と、$\Delta x\,\Delta k$ を求めましょう。運動量のはば $\Delta p$ も求めます。
\phbamen
数学の本（9-4⑨）の「はば $\sigma$ と $1/\sigma$」を、電子の波束に当てはめる場面です（⑦）。
\phmushi
波束が時間とともに広がることは考えず、ある瞬間の形だけを見ます。
\phzu
$|\psi|^2\propto e^{-x^2/(2\Delta x^2)}$、波の高さ $\psi\propto e^{-x^2/(4\Delta x^2)}$。
\phriyu
高さのはばは $\sigma=\sqrt2\,\Delta x$ で、9-4⑨から波数の高さのはばは $1/\sigma$。確率は2乗なので、はばは $1/\sqrt2$ 倍です（⑦）。
\phtateru
\[ \Delta k=\frac{1}{\sigma}\times\frac{1}{\sqrt2}=\frac{1}{\sqrt2\,\Delta x}\times\frac{1}{\sqrt2}=\frac{1}{2\Delta x},\qquad\Delta p=\hbar\,\Delta k \]
\phtoku
\[ \Delta k=\frac{1}{2\times1.0\times10^{-9}}=\qty{5.0e8}{rad/m},\qquad\Delta x\,\Delta k=0.50 \]
\[ \Delta p=1.055\times10^{-34}\times5.0\times10^{8}\approx\qty{5.3e-26}{kg.m/s} \]
\phtashikame
単位：$\Delta k$ は 1/m（rad/m）、$\Delta x\,\Delta k$ は単位なし。\\
大きさ：ちょうど $\tfrac12$ で、不確定性関係の等号です。ガウスの形が、いちばん「むだのない」波束です。実際に、細かく分けて数値で変換しても（\texttt{calc\_p5\_03.py}）、0.50 になりました。\\
極端な場合：$\Delta x$ を10倍にすると、$\Delta k$ は10分の1になります。
\kotae $\Delta k=\qty{5.0e8}{rad/m}$、$\Delta x\,\Delta k=0.50$、$\Delta p\approx\qty{5.3e-26}{kg.m/s}$
\end{phreidai}
\begin{jibun}
例題10 の波束のはばを $\Delta x=\qty{0.50}{nm}$ にすると、$\Delta k$ はいくらですか。
\bkjans{$\Delta k=1/(2\Delta x)=\qty{1.0e9}{rad/m}$（2倍）}
\end{jibun}

\label{ph:原-3:ex11}
\begin{phreidai}[電子の波の位相の速さと群速度]
100 V で加速した電子の、位相の速さと群速度を求め、電子の速さとくらべましょう。
\phbamen
電子の波の山が進む速さと、波束のかたまりが進む速さを分けて考える場面です（⑨）。
\phmushi
相対論は考えません。エネルギーの0は、止まった電子にとります。
\phzu
分散関係 $\omega=\hbar k^2/(2m)$。電子の速さ $v=\qty{5.93e6}{m/s}$（例題2）。
\phriyu
位相の速さは $\omega/k$、群速度は $\dd\omega/\dd k$（⑨・波-5 ⑦）。
\phtateru
\[ \frac{\omega}{k}=\frac{\hbar k}{2m},\qquad v_{\mathrm g}=\frac{\dd\omega}{\dd k}=\frac{\hbar k}{m},\qquad\hbar k=p=mv \]
\phtoku
\[ \frac{\omega}{k}=\frac{v}{2}\approx\qty{2.97e6}{m/s},\qquad v_{\mathrm g}=v\approx\qty{5.93e6}{m/s} \]
\phtashikame
単位：m/s。\\
大きさ：群速度が電子の速さと同じなので、「電子がどこにいるか」（波束のかたまり）は、つぶとしての電子と同じように動きます。\\
極端な場合：光の波（$\omega=ck$）では、位相の速さと群速度がどちらも $c$ で等しくなります。
\kotae 位相の速さ約\qty{3.0e6}{m/s}（電子の速さの半分）、群速度約\qty{5.9e6}{m/s}（電子の速さと同じ）
\end{phreidai}
\begin{jibun}
1.0 kV で加速した電子の、群速度と位相の速さを求めましょう。
\bkjans{$v=\sqrt{2eV/m}\approx\qty{1.9e7}{m/s}$ が群速度。位相の速さはその半分で約\qty{9.4e6}{m/s}}
\end{jibun}

\label{ph:原-3:ex12}
\begin{phreidai}[電子顕微鏡で見分けられる細かさ]
200 kV の透過電子顕微鏡で、レンズの開き角を $\alpha=\qty{10}{mrad}$（0.010 rad）とします。電子の波長（相対論で直す）と、見分けられる細かさ $\delta\approx0.61\lambda/\alpha$ を求めましょう。
\phbamen
はてな？の答えを、数で確かめる場面です（⑥）。
\phmushi
レンズの収差は、使える角 $\alpha$ を小さくすることだけで考えに入れます。試料の中での電子の散らばりや、機械のゆれは無視します。
\phzu
$V=\qty{2.00e5}{V}$、$\alpha=0.010$ rad、$mc^2=\qty{511}{keV}$。
\phriyu
波長は②（相対論の直しつき）。見分けられる細かさは、円い口の回折（波-13 ⑤）の $0.61\lambda/\alpha$。
\phtateru
\[ \lambda=\frac{1.226\ \mathrm{nm}}{\sqrt{V}}\times\frac{1}{\sqrt{1+eV/(2mc^2)}},\qquad\delta=\frac{0.61\lambda}{\alpha} \]
\phtoku
\[ \frac{1.226}{\sqrt{2.00\times10^5}}=\qty{2.74e-3}{nm},\qquad1+\frac{200}{1022}=1.196 \]
\[ \lambda=\frac{\qty{2.74}{pm}}{\sqrt{1.196}}\approx\qty{2.51}{pm} \]
\[ \delta=\frac{0.61\times2.51\times10^{-12}}{0.010}\approx\qty{1.5e-10}{m}=\qty{0.15}{nm} \]
\phtashikame
単位：m÷rad＝m。\\
大きさ：原子1個（約 0.1〜0.3 nm）を見分けられる細かさで、はてな？の答えと合います。光の顕微鏡（約 \qty{0.2}{\micro m}）の約1000分の1です。\\
極端な場合：収差をなくすレンズ（収差補正）で $\alpha$ を大きくできれば、もっと細かく見えます。
\kotae 波長約\qty{2.51}{pm}、約\qty{0.15}{nm}
\end{phreidai}
\begin{jibun}
300 kV の電子顕微鏡で、開き角 0.010 rad のとき、見分けられる細かさはいくらですか。相対論で直した波長は \qty{1.97}{pm} です。
\bkjans{$\delta=0.61\lambda/\alpha$ で、約\qty{0.12}{nm}}
\end{jibun}

%% ======================== まちがえやすいところ ========================
\Needspace{11\baselineskip}\bksection{tsumazuki}{まちがえやすいところをチェック}
\begin{tsumazuki}{電子の波長を $hc/E$ で求める？}
\bkng{54 eV の電子の波長は $1240\div54\approx\qty{23}{nm}$。}
\bkok{$hc/E$ は光子だけの式。電子は $\lambda=h/p$、$p=\sqrt{2mE}$ で \qty{0.167}{nm}。}
光子は $p=E/c$ ですが、電子の運動量は $\sqrt{2mE}$ です。同じエネルギーでも、電子の波長は光子よりずっと短くなります（54 eV で約140分の1）。
\end{tsumazuki}

\begin{tsumazuki}{電圧を2倍にすると、波長は半分？}
\bkng{$V$ を2倍にすると、$\lambda$ は $\tfrac12$ になる。}
\bkok{$\lambda\propto1/\sqrt V$ なので $1/\sqrt2\approx0.71$ 倍。半分にするには4倍。}
運動エネルギーは $V$ に比例しますが、運動量はその平方根に比例するからです（②の2倍にすると）。
\end{tsumazuki}

\begin{tsumazuki}{不確定性は、はかり方が下手だから？}
\bkng{もっとよい測定器を作れば、位置も運動量もぴったり決められる。}
\bkok{電子が波なので、両方が同時にぴったり決まった状態が、もともとない。}
波束は、位置をせまくすると波長が広がる、という波の性質そのものです（⑦）。
\end{tsumazuki}

\begin{tsumazuki}{電子1個が、半分ずつに分かれて通る？}
\bkng{電子が2つのすきまで半分に分かれ、スクリーンで2つの点になる。}
\bkok{スクリーンでは、いつも電子1個の1つの点。届く場所の確率が干渉のしまの形になる。}
半分の電荷の電子が見つかったことはありません。干渉するのは、確率を決める波 $\psi$ です（⑤）。
\end{tsumazuki}

%% ======================== 工学ではこう使う ========================
\Needspace{18\baselineskip}\bksection{phgear}{工学ではこう使う}
\begin{phkougaku}[電子の波で、材料と部品を見る]
\textbf{半導体の工場}：スマホのチップのトランジスタは、いちばん細い所が数 nm から十数 nm しかありません。工場では、作ったチップの一部を切り取って、透過電子顕微鏡で断面の形を見たり、走査電子顕微鏡で表面の線のはばをはかったりして、正しく作れているかを確かめます。

\textbf{材料の研究}：鉄やアルミニウムの強さは、中の結晶のつぶの大きさや、原子の並びの乱れで決まります。電子顕微鏡で原子の並びを見て、電子線回折（④）で結晶の種類と向きを調べます。新しい電池の材料や、触媒（化学反応を助ける物）のつぶも、こうして調べます。

\textbf{生き物を見る}：たんぱく質の分子を水ごとすばやく凍らせ、電子顕微鏡でたくさんの向きから写して、分子の立体の形を組み立てる方法（クライオ電子顕微鏡）が使われています。薬の開発に役立っています。

\begin{phketa}[電子の波の目安]
\phketaline{電子の波長}{100 V で 0.12 nm、10 kV で 12 pm、100 kV で 3.7 pm}
\phketaline{見分けられる細かさ}{光の顕微鏡 約\qty{0.2}{\micro m}、SEM 約 1 nm、TEM 約 0.1 nm}
\phketaline{不確定性の目安}{原子の大きさ 0.1 nm に閉じこめると約 1 eV}
\phboxhead{答えが桁で合っているか}
電子の波長は、$1.226/\sqrt V$ nm でざっと確かめます。kV と V、pm と nm の直し忘れで、1000倍間違えやすいので気をつけましょう。数十 kV より上は、相対論の直し（数\,\%）を忘れないこと。
\end{phketa}
\end{phkougaku}
\phstation{量子力学,電子顕微鏡学,材料評価,半導体工学}{ド・ブロイ波長、電子線回折、透過電子顕微鏡（TEM）、走査電子顕微鏡（SEM）、波束、不確定性関係}

%% ======================== 章末問題 ========================
\bkendproblems
\bkproblemgroup{基本}{3}
\bkproblem 900 V で加速した電子のド・ブロイ波長を求めましょう。100 V のときの何倍ですか。
\bkpans{$1.226\div\sqrt{900}\approx\qty{0.041}{nm}$。100 V（\qty{0.123}{nm}）の $\tfrac13$ 倍}
\bkproblem 25 eV の電子と、25 eV の光子の波長を、それぞれ求めましょう。
\bkpans{電子 $1.226\div\sqrt{25}\approx\qty{0.25}{nm}$、光子 $1240\div25\approx\qty{50}{nm}$}
\bkproblem 波長が \qty{0.20}{nm} の中性子（質量 \qty{1.675e-27}{kg}）の速さを求めましょう。
\bkpans{$v=h/(m_{\mathrm n}\lambda)$ で、約\qty{2.0e3}{m/s}}
\bkproblem 波長 \qty{5.0}{pm} の電子を、間かく \qty{2.0}{\micro m} の2つのすきまに通し、\qty{1.5}{m} 先のスクリーンに当てます。しまの間かくを求めましょう。
\bkpans{$\Delta x=L\lambda/d$ で、約\qty{3.8}{\micro m}}
\bkproblem ニッケルの表面（列の間かく \qty{0.215}{nm}）に垂直に電子を当てると、40° の向きに $n=1$ の強め合いが見えました。電子の波長と、加速電圧を求めましょう。
\bkpans{$\lambda=0.215\times\sin\ang{40}\approx\qty{0.138}{nm}$。$V=(1.226/0.138)^2\approx\qty{79}{V}$}
\bkproblemgroup{チャレンジ}{4}
\bkproblem 陽子を、原子核くらいの $\Delta x=\qty{2.0}{fm}$（\qty{2.0e-15}{m}）の中に閉じこめると、運動量のばらつきの最小の値と、運動エネルギーの目安 $(\Delta p)^2/(2m_{\mathrm p})$ はいくらですか（陽子の質量 \qty{1.67e-27}{kg}）。原子核の中のエネルギーが MeV の大きさになるわけを説明しましょう。
\bkpans{$\Delta p=\hbar/(2\Delta x)$ で約\qty{2.6e-20}{kg.m/s}。運動エネルギーの目安は約\qty{2.1e-13}{J}、約\qty{1.3}{MeV}。原子の約10万分の1のせまい所に閉じこめるので、運動エネルギーが大きくなる}
\bkproblem 4.0 kV で加速した電子の、ド・ブロイ波長・群速度・位相の速さを求めましょう（相対論は考えない）。
\bkpans{$1.226\div\sqrt{4000}\approx\qty{0.019}{nm}$（約\qty{19}{pm}）、群速度は電子の速さ $\sqrt{2eV/m}\approx\qty{3.8e7}{m/s}$、位相の速さはその半分の約\qty{1.9e7}{m/s}}
\bkproblem ある原子の段の寿命が \qty{10}{ns} です。この段から出る光の自然幅は何 MHz ですか。また、寿命を10倍にすると、はばはどうなりますか。
\bkpans{$\Delta f=1/(2\pi\tau)$ で約\qty{16}{MHz}。寿命を10倍にすると、はばは $\tfrac{1}{10}$（約\qty{1.6}{MHz}）}

\begin{furikaeri}
\begin{checklist}
\item 加速した電圧から、電子の波長（ド・ブロイ波長）を求められる。\bkfurref{ph:原-3:sec2}
\item 電子の回折と、電子1個ずつの二重スリットのしま、電子顕微鏡が細かい物を見分けられるわけを、波の式で説明できる。\bkfurref{ph:原-3:sec5}
\item 波の集まり（波束）のはばから不確定性関係を導き、原子の中の電子のエネルギーの大きさを見積もれる。\bkfurref{ph:原-3:sec7}
\item 電子の波の分散関係から、群速度が粒の速さになることを示せる。\bkfurref{ph:原-3:sec9}
\end{checklist}
\end{furikaeri}
\bkfukidashi{sora}{電子は、1個ずつ見ればつぶ、たくさん集めれば波のしま。波長がとても短いから、電子顕微鏡は原子まで見えるんだね。位置と速さが同時に決まらないのも、波だからなんだ！}
